\documentclass[a4paper,10pt]{article}
\usepackage{eumat}

\begin{document}
\begin{eulernotebook}
\eulerheading{EMT untuk Perhitungan Aljabar (Bagian 1)}
\begin{eulercomment}
Pada notebook ini Anda belajar menggunakan EMT untuk melakukan berbagai
perhitungan terkait materi Aljabar. Setiap bagian berisi:

- Uraian materi (teori);
- Contoh perhitungan dengan EMT;
- Soal latihan beserta penyelesaiannya di EMT;
- Penjelasan hasil.

Cara memakai: pindahkan kursor ke baris perintah (diawali tanda "\textgreater{}"),
lalu tekan ENTER. Keluaran yang belum tampil akan muncul setelah dijalankan.

\end{eulercomment}
\eulersubheading{Contoh pertama}
\begin{eulercomment}
Uraian materi: sifat-sifat eksponen untuk a, b tak nol dan bilangan bulat m, n:

\end{eulercomment}
\begin{eulerformula}
\[
a^m \cdot a^n = a^{m+n}, \quad \frac{a^m}{a^n}=a^{m-n}, \quad (a^m)^n=a^{mn}, \quad a^{-n}=\frac{1}{a^n}
\]
\end{eulerformula}
\begin{eulercomment}
Menyederhanakan bentuk aljabar:

\end{eulercomment}
\begin{eulerformula}
\[
6x^{-3}y^5\times -7x^2y^{-9}
\]
\end{eulerformula}
\begin{eulercomment}
\end{eulercomment}
\begin{eulerprompt}
>$&6*x^(-3)*y^5*-7*x^2*y^(-9)
\end{eulerprompt}
\begin{eulercomment}
Menjabarkan:

\end{eulercomment}
\begin{eulerformula}
\[
(6x^{-3}+y^5)(-7x^2-y^{-9})
\]
\end{eulerformula}
\begin{eulerprompt}
>$&showev('expand((6*x^(-3)+y^5)*(-7*x^2-y^(-9))))
\end{eulerprompt}
\begin{eulercomment}
Penjelasan: koefisien dikalikan (6 x -7 = -42), pangkat variabel yang sama
dijumlahkan (x\textasciicircum{}(-3+2)=x\textasciicircum{}(-1), y\textasciicircum{}(5-9)=y\textasciicircum{}(-4)). Hasilnya -42/(x y\textasciicircum{}4).

*** Latihan 1.1
Sederhanakan (3a\textasciicircum{}2 b\textasciicircum{}-1)(4a\textasciicircum{}-5 b\textasciicircum{}3). Hasil yang diharapkan: 12 b\textasciicircum{}2/a\textasciicircum{}3.

\end{eulercomment}
\begin{eulerformula}
\[
(3a^{2}b^{-1})(4a^{-5}b^{3})
\]
\end{eulerformula}
\begin{eulerprompt}
>$&(3*a^2*b^(-1))*(4*a^(-5)*b^3)
\end{eulerprompt}
\begin{eulercomment}
*** Latihan 1.2
Jabarkan (2x-3)(x\textasciicircum{}2+x-5). Hasil yang diharapkan: 2x\textasciicircum{}3 - x\textasciicircum{}2 - 13x + 15.

\end{eulercomment}
\begin{eulerformula}
\[
(2x-3)(x^2+x-5)
\]
\end{eulerformula}
\begin{eulerprompt}
>$&expand((2*x-3)*(x^2+x-5))
\end{eulerprompt}
\begin{eulercomment}
*** Latihan 1.3
Faktorkan x\textasciicircum{}3-6x\textasciicircum{}2+11x-6. Hasil yang diharapkan: (x-1)(x-2)(x-3).

\end{eulercomment}
\begin{eulerformula}
\[
x^3-6x^2+11x-6
\]
\end{eulerformula}
\begin{eulerprompt}
>$&factor(x^3-6*x^2+11*x-6)
\end{eulerprompt}
\begin{eulercomment}
Penjelasan: expand() menjabarkan perkalian suku banyak, sedangkan factor()
melakukan kebalikannya. Hasil factor() dapat dicek dengan expand().
\end{eulercomment}
\begin{eulerprompt}
>$&expand(factor(x^3-6*x^2+11*x-6))
\end{eulerprompt}
\begin{eulercomment}
\end{eulercomment}
\eulersubheading{Baris Perintah (The Command Line)}
\begin{eulercomment}
Uraian materi: sebuah baris perintah Euler terdiri atas satu atau beberapa
perintah yang diakhiri titik koma ";" atau koma ",". Titik koma mencegah
hasil dicetak. Koma setelah perintah terakhir boleh dihilangkan.

Baris perintah berikut hanya mencetak hasil ekspresi, bukan penugasan
maupun perintah format.
\end{eulercomment}
\begin{eulerprompt}
>r:=2; h:=4; pi*r^2*h/3
\end{eulerprompt}
\begin{euleroutput}
  16.7551608191
\end{euleroutput}
\begin{eulercomment}
Perintah harus dipisahkan dengan spasi. Baris berikut mencetak dua hasil.
\end{eulercomment}
\begin{eulerprompt}
>pi*2*r*h, %+2*pi*r*h // Ingat tanda % menyatakan hasil perhitungan terakhir sebelumnya
\end{eulerprompt}
\begin{euleroutput}
  50.2654824574
  100.530964915
\end{euleroutput}
\begin{eulercomment}
Baris perintah dijalankan sesuai urutan saat pengguna menekan return.
Jadi Anda mendapat nilai baru setiap kali baris kedua dijalankan.
\end{eulercomment}
\begin{eulerprompt}
>x := 1;
>x := cos(x) // nilai cosinus (x dalam radian)
\end{eulerprompt}
\begin{euleroutput}
  0.540302305868
\end{euleroutput}
\begin{eulerprompt}
>x := cos(x)
\end{eulerprompt}
\begin{euleroutput}
  0.857553215846
\end{euleroutput}
\begin{eulercomment}
Jika dua baris dihubungkan dengan "...", kedua baris selalu dijalankan
bersamaan.
\end{eulercomment}
\begin{eulerprompt}
>x := 1.5; ...
>x := (x+2/x)/2, x := (x+2/x)/2, x := (x+2/x)/2,
\end{eulerprompt}
\begin{euleroutput}
  1.41666666667
  1.41421568627
  1.41421356237
\end{euleroutput}
\begin{eulercomment}
Ini juga cara yang baik untuk memecah perintah panjang menjadi beberapa
baris. Tekan Ctrl+Return untuk memecah baris pada posisi kursor, dan
Ctrl+Back untuk menggabungkan baris.

Untuk melipat semua multi-baris tekan Ctrl+L. Untuk melipat satu multi-baris,
awali baris pertamanya dengan "\%+ ".
\end{eulercomment}
\begin{eulerprompt}
>%+ x=4+5; ...
\end{eulerprompt}
\begin{eulercomment}
Baris yang diawali \%\% tidak akan ditampilkan sama sekali.
\end{eulercomment}
\begin{euleroutput}
  81
\end{euleroutput}
\begin{eulercomment}
Euler mendukung perulangan pada baris perintah, selama muat dalam satu baris
atau satu multi-baris. Di dalam program, batasan ini tentu tidak berlaku.
\end{eulercomment}
\begin{eulerprompt}
>x=1; for i=1 to 5; x := (x+2/x)/2, end; // menghitung akar 2
\end{eulerprompt}
\begin{euleroutput}
  1.5
  1.41666666667
  1.41421568627
  1.41421356237
  1.41421356237
\end{euleroutput}
\begin{eulercomment}
Multi-baris diperbolehkan. Pastikan baris diakhiri " ...".
\end{eulercomment}
\begin{eulerprompt}
>x := 1.5; // komentar ditulis sebelum tanda ...
>repeat xnew:=(x+2/x)/2; until xnew~=x; ...
>   x := xnew; ...
>end; ...
>x,
\end{eulerprompt}
\begin{euleroutput}
  1.41421356237
\end{euleroutput}
\begin{eulercomment}
Struktur kondisional juga dapat digunakan.
\end{eulercomment}
\begin{eulerprompt}
>if E^pi>pi^E; then "Thought so!", endif;
\end{eulerprompt}
\begin{euleroutput}
  Thought so!
\end{euleroutput}
\begin{eulercomment}
Saat menjalankan perintah, kursor boleh berada di posisi mana pun pada baris
perintah. Gunakan tombol panah untuk berpindah antarperintah. Saat kursor
bergerak, pasangan kurung buka dan tutup akan tersorot. Perhatikan juga
baris status: setelah kurung buka fungsi sqrt(), baris status menampilkan
teks bantuan fungsi tersebut.
\end{eulercomment}
\begin{eulerprompt}
>sqrt(sin(10°)/cos(20°))
\end{eulerprompt}
\begin{euleroutput}
  Variable  not found!
  Error in degree
  Error in:
  sqrt(sin(10°)/cos(20°)) ...
               ^
\end{euleroutput}
\begin{eulercomment}
Untuk melihat bantuan perintah terakhir, buka jendela bantuan dengan F1.
Klik ganda pada perintah juga membuka bantuannya. Coba klik ganda "exp".
\end{eulercomment}
\begin{eulerprompt}
>exp(log(2.5))
\end{eulerprompt}
\begin{euleroutput}
  2.5
\end{euleroutput}
\begin{eulercomment}
Salin dan tempel juga dapat dilakukan dengan Ctrl-C dan Ctrl-V.
\end{eulercomment}
\begin{eulercomment}

Uraian materi (metode Newton/Heron untuk akar kuadrat): barisan

\end{eulercomment}
\begin{eulerformula}
\[
x_{n+1}=\frac{1}{2}\left(x_n+\frac{a}{x_n}\right)
\]
\end{eulerformula}
\begin{eulercomment}
konvergen ke akar kuadrat dari a untuk x\_0 \textgreater{} 0.

*** Latihan 2.1
Hitung volume kerucut dengan jari-jari 3 dan tinggi 5.
Rumus: V = (1/3) pi r\textasciicircum{}2 h. Hasil yang diharapkan: 15 pi, sekitar 47.1239.
\end{eulercomment}
\begin{eulerprompt}
>r:=3; h:=5; pi*r^2*h/3
\end{eulerprompt}
\begin{euleroutput}
  47.1238898038
\end{euleroutput}
\begin{eulerprompt}
>$&15*%pi, &float(%)   // pembanding eksak: 15 pi
\end{eulerprompt}
\begin{euleroutput}
  
                            47.12388980384689
  
\end{euleroutput}
\begin{eulercomment}
*** Latihan 2.2
Hitung akar 3 dengan iterasi x -\textgreater{} (x+3/x)/2 mulai dari x=2 sebanyak 5 kali.
Barisan yang diharapkan: 1.75, 1.7321, 1.73205, ... mendekati sqrt(3).
\end{eulercomment}
\begin{eulerprompt}
>x=2; for i=1 to 5; x := (x+3/x)/2, end; // menghitung akar 3
\end{eulerprompt}
\begin{euleroutput}
  1.75
  1.73214285714
  1.73205081001
  1.73205080757
  1.73205080757
\end{euleroutput}
\begin{eulerprompt}
>sqrt(3)  // pembanding
\end{eulerprompt}
\begin{euleroutput}
  1.73205080757
\end{euleroutput}
\begin{eulercomment}
Penjelasan: setiap iterasi menggandakan jumlah digit yang benar (konvergen
kuadratik), sehingga setelah sekitar 4 iterasi nilainya sudah sama dengan
sqrt(3) pada ketelitian 12 digit.

\end{eulercomment}
\eulersubheading{Sintaks Dasar}
\begin{eulercomment}
Euler mengenal fungsi-fungsi matematika yang umum. Fungsi trigonometri
bekerja dalam radian atau derajat. Untuk mengubah ke derajat, tambahkan
simbol derajat (tombol F7) pada nilai, atau gunakan fungsi rad(x). Fungsi
akar kuadrat bernama sqrt. Tentu x\textasciicircum{}(1/2) juga dapat digunakan.

Untuk memberi nilai variabel, gunakan "=" atau ":=". Demi kejelasan,
pengantar ini memakai bentuk kedua. Spasi tidak berpengaruh, tetapi spasi
antarperintah diperlukan.

Beberapa perintah dalam satu baris dipisahkan dengan "," atau ";". Titik
koma menekan keluaran. Di akhir baris, "," dianggap ada jika ";" tidak ada.
\end{eulercomment}
\begin{eulerprompt}
>g:=9.81; t:=2.5; 1/2*g*t^2
\end{eulerprompt}
\begin{euleroutput}
  30.65625
\end{euleroutput}
\begin{eulercomment}
EMT memakai sintaks pemrograman untuk ekspresi. Untuk memasukkan

\end{eulercomment}
\begin{eulerformula}
\[
e^2 \cdot \left( \frac{1}{3+4 \log(0.6)}+\frac{1}{7} \right)
\]
\end{eulerformula}
\begin{eulercomment}
Anda harus memasang kurung dengan benar dan memakai / untuk pecahan.
Konstanta Euler e dituliskan E di EMT.
\end{eulercomment}
\begin{eulerprompt}
>E^2*(1/(3+4*log(0.6))+1/7)
\end{eulerprompt}
\begin{euleroutput}
  8.77908249441
\end{euleroutput}
\begin{eulercomment}
Untuk menghitung ekspresi rumit seperti

\end{eulercomment}
\begin{eulerformula}
\[
\left(\frac{\frac17 + \frac18 + 2}{\frac13 + \frac12}\right)^2 \pi
\]
\end{eulerformula}
\begin{eulercomment}
tuliskan dalam bentuk satu baris.
\end{eulercomment}
\begin{eulerprompt}
>((1/7 + 1/8 + 2) / (1/3 + 1/2))^2 * pi
\end{eulerprompt}
\begin{euleroutput}
  23.2671801626
\end{euleroutput}
\begin{eulercomment}
Pasang kurung dengan hati-hati pada sub-ekspresi yang harus dihitung lebih
dulu. Huruf Yunani pi harus dituliskan "pi".

Hasil perhitungan ini bilangan floating point, yang secara default dicetak
dengan sekitar 12 digit. Berikut cara merujuk hasil sebelumnya dalam satu
baris yang sama.
\end{eulercomment}
\begin{eulerprompt}
>1/3+1/7, fraction %
\end{eulerprompt}
\begin{euleroutput}
  0.47619047619
  10/21
\end{euleroutput}
\begin{eulercomment}
Perintah Euler dapat berupa ekspresi atau perintah primitif. Ekspresi
tersusun dari operator dan fungsi. Jika perlu, beri kurung untuk memaksa
urutan pengerjaan yang benar. Jika ragu, memasang kurung adalah ide baik.
\end{eulercomment}
\begin{eulerprompt}
>(cos(pi/4)+1)^3*(sin(pi/4)+1)^2
\end{eulerprompt}
\begin{euleroutput}
  14.4978445072
\end{euleroutput}
\begin{eulercomment}
Operator numerik Euler antara lain:

\end{eulercomment}
\begin{eulerttcomment}
 + operator plus (uner atau biner)
 - operator minus (uner atau biner)
 *, /
 . perkalian matriks
 a^b pangkat untuk a positif atau b bulat (a**b juga bisa)
 n! operator faktorial
\end{eulerttcomment}
\begin{eulercomment}

dan masih banyak lagi.

Beberapa fungsi yang mungkin Anda perlukan:

\end{eulercomment}
\begin{eulerttcomment}
 sin,cos,tan,atan,asin,acos,rad,deg
 log,exp,log10,sqrt,logbase
 bin,logbin,logfac,mod,floor,ceil,round,abs,sign
 conj,re,im,arg,conj,real,complex
 beta,betai,gamma,complexgamma,ellrf,ellf,ellrd,elle
 bitand,bitor,bitxor,bitnot
\end{eulerttcomment}
\begin{eulercomment}

Beberapa perintah punya alias, misalnya ln untuk log.
\end{eulercomment}
\begin{eulerprompt}
>ln(E^2), arctan(tan(0.5))
\end{eulerprompt}
\begin{euleroutput}
  2
  0.5
\end{euleroutput}
\begin{eulerprompt}
>sin(30°)
\end{eulerprompt}
\begin{euleroutput}
  Variable  not found!
  Error in degree
  Error in:
  sin(30°) ...
          ^
\end{euleroutput}
\begin{eulercomment}
Gunakan kurung setiap kali ragu pada urutan pengerjaan! Perintah berikut
tidak sama dengan (2\textasciicircum{}3)\textasciicircum{}4, yang merupakan default untuk 2\textasciicircum{}3\textasciicircum{}4 di EMT
(sebagian sistem numerik melakukannya sebaliknya).
\end{eulercomment}
\begin{eulerprompt}
>2^3^4, (2^3)^4, 2^(3^4)
\end{eulerprompt}
\begin{euleroutput}
  2.41785163923e+24
  4096
  2.41785163923e+24
\end{euleroutput}
\begin{eulercomment}
Uraian materi (urutan operasi): pangkat dikerjakan lebih dulu, lalu
perkalian/pembagian, lalu penjumlahan/pengurangan; kurung didahulukan.

*** Latihan 3.1
Hitung (2/3 + 3/4) / (5/6 - 1/2) dalam bentuk pecahan.
Hasil yang diharapkan: 17/4 (= 4.25).
\end{eulercomment}
\begin{eulerprompt}
>fraction (2/3+3/4)/(5/6-1/2)
\end{eulerprompt}
\begin{euleroutput}
  17/4
\end{euleroutput}
\begin{eulercomment}
*** Latihan 3.2
Hitung e\textasciicircum{}2 * (1/(3+4 ln 0,6) + 1/7) dan bandingkan dengan contoh di atas
(harus sama: 8.77908249441).
\end{eulercomment}
\begin{eulerprompt}
>E^2*(1/(3+4*ln(0.6))+1/7)
\end{eulerprompt}
\begin{euleroutput}
  8.77908249441
\end{euleroutput}
\begin{eulercomment}
*** Latihan 3.3
Hitung sin(30°)\textasciicircum{}2 + cos(30°)\textasciicircum{}2. Hasil yang diharapkan: 1 (identitas
Pythagoras sin\textasciicircum{}2 + cos\textasciicircum{}2 = 1).
\end{eulercomment}
\begin{eulerprompt}
>sin(30°)^2+cos(30°)^2
\end{eulerprompt}
\begin{euleroutput}
  Variable  not found!
  Error in degree
  Error in:
  sin(30°)^2+cos(30°)^2 ...
          ^
\end{euleroutput}
\begin{eulercomment}
Penjelasan: fraction mengubah hasil desimal menjadi pecahan terdekat. Pada
latihan 3.3, hasilnya 1 sampai ketelitian mesin; sisa galat sangat kecil
tidak ditampilkan pada format default.

\end{eulercomment}
\eulersubheading{Bilangan Real}
\begin{eulercomment}
Tipe data utama di Euler adalah bilangan real. Bilangan real disimpan dalam
format IEEE dengan ketelitian sekitar 16 digit desimal.
\end{eulercomment}
\begin{eulerprompt}
>longest 1/3
\end{eulerprompt}
\begin{euleroutput}
       0.3333333333333333 
\end{euleroutput}
\begin{eulercomment}
Representasi biner internalnya memakai 8 byte.
\end{eulercomment}
\begin{eulerprompt}
>printdual(1/3)
\end{eulerprompt}
\begin{euleroutput}
  1.0101010101010101010101010101010101010101010101010101*2^-2
\end{euleroutput}
\begin{eulerprompt}
>printhex(1/3)
\end{eulerprompt}
\begin{euleroutput}
  5.5555555555554*16^-1
\end{euleroutput}
\begin{eulercomment}
*** Latihan 4.1
Tampilkan sqrt(2) dengan semua digit valid, lalu kuadratkan kembali.
Perhatikan selisihnya dari 2 (galat pembulatan).
\end{eulercomment}
\begin{eulerprompt}
>longest sqrt(2)
\end{eulerprompt}
\begin{euleroutput}
        1.414213562373095 
\end{euleroutput}
\begin{eulerprompt}
>longest sqrt(2)^2-2
\end{eulerprompt}
\begin{euleroutput}
    4.440892098500626e-16 
\end{euleroutput}
\begin{eulercomment}
Penjelasan: sqrt(2) tidak dapat dinyatakan persis dalam biner, sehingga
sqrt(2)\textasciicircum{}2 tidak tepat 2. Selisihnya sekitar 4.4e-16, di orde epsilon mesin.

\end{eulercomment}
\eulersubheading{String}
\begin{eulercomment}
String di Euler didefinisikan dengan "...".
\end{eulercomment}
\begin{eulerprompt}
>"A string can contain anything."
\end{eulerprompt}
\begin{euleroutput}
  A string can contain anything.
\end{euleroutput}
\begin{eulercomment}
String dapat digabung dengan \textbar{} atau +. Ini juga berlaku untuk bilangan,
yang akan diubah menjadi string.
\end{eulercomment}
\begin{eulerprompt}
>"The area of the circle with radius " + 2 + " cm is " + pi*4 + " cm^2."
\end{eulerprompt}
\begin{euleroutput}
  The area of the circle with radius 2 cm is 12.5663706144 cm^2.
\end{euleroutput}
\begin{eulercomment}
Fungsi print juga mengubah bilangan menjadi string. Fungsi ini menerima
jumlah digit dan jumlah tempat desimal (0 untuk keluaran padat).
\end{eulercomment}
\begin{eulerprompt}
>"Golden Ratio : " + print((1+sqrt(5))/2,5,0)
\end{eulerprompt}
\begin{euleroutput}
  Golden Ratio : 1.61803
\end{euleroutput}
\begin{eulercomment}
Ada string khusus none yang tidak dicetak.
\end{eulercomment}
\begin{eulerprompt}
>none
\end{eulerprompt}
\begin{eulercomment}
Untuk mengubah string menjadi bilangan, cukup evaluasi.
\end{eulercomment}
\begin{eulerprompt}
>"1234.5"()
\end{eulerprompt}
\begin{euleroutput}
  1234.5
\end{euleroutput}
\begin{eulercomment}
Vektor string dibuat dengan notasi vektor [...].
\end{eulercomment}
\begin{eulerprompt}
>v:=["affe","charlie","bravo"]
\end{eulerprompt}
\begin{euleroutput}
  affe
  charlie
  bravo
\end{euleroutput}
\begin{eulercomment}
Vektor string kosong dinotasikan [none]. Vektor string dapat digabung.
\end{eulercomment}
\begin{eulerprompt}
>w:=[none]; w|v|v
\end{eulerprompt}
\begin{euleroutput}
  affe
  charlie
  bravo
  affe
  charlie
  bravo
\end{euleroutput}
\begin{eulercomment}
String dapat memuat karakter Unicode (kode UTF-8). Gunakan u"..." dan salah
satu entitas HTML. Keluaran bisa tampak rusak di sebagian tampilan
(masalah penyandian, bukan kesalahan perintah).
\end{eulercomment}
\begin{eulerprompt}
>u"&alpha; = " + 45 + u"&deg;" // pdfLaTeX mungkin gagal menampilkan secara benar
\end{eulerprompt}
\begin{euleroutput}
  α = 45°
\end{euleroutput}
\begin{eulercomment}
Fungsi strtochar() mengenali string Unicode dan menerjemahkannya dengan
benar. Hasilnya vektor bilangan Unicode. Fungsi kebalikannya chartoutf().
\end{eulercomment}
\begin{eulerprompt}
>v=strtochar(u"&Auml; is a German letter")
\end{eulerprompt}
\begin{euleroutput}
  [196,  32,  105,  115,  32,  97,  32,  71,  101,  114,  109,  97,  110,
  32,  108,  101,  116,  116,  101,  114]
\end{euleroutput}
\begin{eulerprompt}
>v[1]=strtochar(u"&Uuml;")[1]; chartoutf(v)
\end{eulerprompt}
\begin{euleroutput}
  Ü is a German letter
\end{euleroutput}
\begin{eulercomment}
Fungsi utf() mengubah string beserta entitasnya menjadi string Unicode.
\end{eulercomment}
\begin{eulerprompt}
>s="We have &alpha;=&beta;."; utf(s) // pdfLaTeX mungkin gagal menampilkan secara benar
\end{eulerprompt}
\begin{euleroutput}
  We have α=β.
\end{euleroutput}
\begin{eulercomment}
Entitas numerik juga dapat dipakai.
\end{eulercomment}
\begin{eulerprompt}
>u"&#196;hnliches"
\end{eulerprompt}
\begin{euleroutput}
  Ähnliches
\end{euleroutput}
\begin{eulercomment}
*** Latihan 5.1
Susun kalimat: "Luas persegi dengan sisi 7 adalah 49" menggunakan
penggabungan string dan angka.
\end{eulercomment}
\begin{eulerprompt}
>"Luas persegi dengan sisi " + 7 + " adalah " + 7^2
\end{eulerprompt}
\begin{euleroutput}
  Luas persegi dengan sisi 7 adalah 49
\end{euleroutput}
\begin{eulercomment}
*** Latihan 5.2
Ubah string "3.5*4" menjadi nilai. Hasil yang diharapkan: 14.
\end{eulercomment}
\begin{eulerprompt}
>"3.5*4"()
\end{eulerprompt}
\begin{euleroutput}
  14
\end{euleroutput}
\begin{eulercomment}
Penjelasan: operator + pada string mengubah angka menjadi teks, sedangkan
memberi tanda kurung () pada string mengevaluasinya sebagai ekspresi.

\end{eulercomment}
\eulersubheading{Nilai Boolean}
\begin{eulercomment}
Nilai Boolean dinyatakan 1=benar dan 0=salah. String dapat dibandingkan,
seperti bilangan.
\end{eulercomment}
\begin{eulerprompt}
>2<1, "apel"<"banana"
\end{eulerprompt}
\begin{euleroutput}
  0
  1
\end{euleroutput}
\begin{eulercomment}
"dan" ditulis "\&\&" dan "atau" ditulis "\textbar{}\textbar{}", seperti bahasa C. (Kata "and"
dan "or" hanya dapat dipakai pada kondisi "if".)
\end{eulercomment}
\begin{eulerprompt}
>2<E && E<3
\end{eulerprompt}
\begin{euleroutput}
  1
\end{euleroutput}
\begin{eulercomment}
Operator Boolean mengikuti aturan bahasa matriks.
\end{eulercomment}
\begin{eulerprompt}
>(1:10)>5, nonzeros(%)
\end{eulerprompt}
\begin{euleroutput}
  [0,  0,  0,  0,  0,  1,  1,  1,  1,  1]
  [6,  7,  8,  9,  10]
\end{euleroutput}
\begin{eulercomment}
Fungsi nonzeros() dapat dipakai untuk memilih elemen tertentu dari vektor.
Pada contoh berikut dipakai kondisi isprime(n).
\end{eulercomment}
\begin{eulerprompt}
>N=2|3:2:99 // N berisi elemen 2 dan bilangan2 ganjil dari 3 s.d. 99
\end{eulerprompt}
\begin{euleroutput}
  [2,  3,  5,  7,  9,  11,  13,  15,  17,  19,  21,  23,  25,  27,  29,
  31,  33,  35,  37,  39,  41,  43,  45,  47,  49,  51,  53,  55,  57,
  59,  61,  63,  65,  67,  69,  71,  73,  75,  77,  79,  81,  83,  85,
  87,  89,  91,  93,  95,  97,  99]
\end{euleroutput}
\begin{eulerprompt}
>N[nonzeros(isprime(N))] //pilih anggota2 N yang prima
\end{eulerprompt}
\begin{euleroutput}
  [2,  3,  5,  7,  11,  13,  17,  19,  23,  29,  31,  37,  41,  43,  47,
  53,  59,  61,  67,  71,  73,  79,  83,  89,  97]
\end{euleroutput}
\begin{eulercomment}
Uraian materi: bilangan prima adalah bilangan asli lebih dari 1 yang hanya
habis dibagi 1 dan dirinya sendiri.

*** Latihan 6.1
Tentukan semua bilangan prima dari 1 sampai 50.
Hasil yang diharapkan: 2,3,5,7,11,13,17,19,23,29,31,37,41,43,47 (15 bilangan).
\end{eulercomment}
\begin{eulerprompt}
>t=1:50; t[nonzeros(isprime(t))]
\end{eulerprompt}
\begin{euleroutput}
  [2,  3,  5,  7,  11,  13,  17,  19,  23,  29,  31,  37,  41,  43,  47]
\end{euleroutput}
\begin{eulerprompt}
>length(nonzeros(isprime(1:50)))
\end{eulerprompt}
\begin{euleroutput}
  15
\end{euleroutput}
\begin{eulercomment}
*** Latihan 6.2
Tentukan semua kelipatan 7 dari 1 sampai 60.
Hasil yang diharapkan: 7,14,21,28,35,42,49,56.
\end{eulercomment}
\begin{eulerprompt}
>t=1:60; t[nonzeros(mod(t,7)==0)]
\end{eulerprompt}
\begin{euleroutput}
  [7,  14,  21,  28,  35,  42,  49,  56]
\end{euleroutput}
\begin{eulercomment}
*** Latihan 6.3
Berapa banyak bilangan dari 1 sampai 100 yang habis dibagi 3 atau 5?
Hasil yang diharapkan: 33 + 20 - 6 = 47 (prinsip inklusi-eksklusi).
\end{eulercomment}
\begin{eulerprompt}
>t=1:100; sum(mod(t,3)==0 || mod(t,5)==0)
\end{eulerprompt}
\begin{euleroutput}
  47
\end{euleroutput}
\begin{eulercomment}
Penjelasan: perbandingan pada vektor menghasilkan vektor 0/1; sum() menghitung
banyaknya elemen yang memenuhi syarat, nonzeros() mengambil indeksnya.

\end{eulercomment}
\eulersubheading{Format Keluaran}
\begin{eulercomment}
Format keluaran default EMT mencetak 12 digit. Agar yakin memakai default,
format direset terlebih dahulu.
\end{eulercomment}
\begin{eulerprompt}
>defformat; pi
\end{eulerprompt}
\begin{euleroutput}
  3.14159265359
\end{euleroutput}
\begin{eulercomment}
Secara internal EMT memakai bilangan double IEEE dengan sekitar 16 digit.
Untuk melihat semua digit, gunakan "longestformat" atau operator "longest".
\end{eulercomment}
\begin{eulerprompt}
>longest pi
\end{eulerprompt}
\begin{euleroutput}
        3.141592653589793 
\end{euleroutput}
\begin{eulercomment}
Representasi heksadesimal internal sebuah bilangan double:
\end{eulercomment}
\begin{eulerprompt}
>printhex(pi)
\end{eulerprompt}
\begin{euleroutput}
  3.243F6A8885A30*16^0
\end{euleroutput}
\begin{eulercomment}
Format keluaran dapat diubah permanen dengan perintah format.
\end{eulercomment}
\begin{eulerprompt}
>format(12,5); 1/3, pi, sin(1)
\end{eulerprompt}
\begin{euleroutput}
      0.33333 
      3.14159 
      0.84147 
\end{euleroutput}
\begin{eulercomment}
Defaultnya adalah format(12).
\end{eulercomment}
\begin{eulerprompt}
>format(12); 1/3
\end{eulerprompt}
\begin{euleroutput}
       0.333333333333 
\end{euleroutput}
\begin{eulercomment}
Fungsi "shortestformat", "shortformat", "longformat" bekerja pada vektor
dengan cara berikut.
\end{eulercomment}
\begin{eulerprompt}
>shortestformat; random(3,8)
\end{eulerprompt}
\begin{euleroutput}
    0.66    0.2   0.89   0.28   0.53   0.31   0.44    0.3 
    0.28   0.88   0.27    0.7   0.22   0.45   0.31   0.91 
    0.19   0.46  0.095    0.6   0.43   0.73   0.47   0.32 
\end{euleroutput}
\begin{eulercomment}
Format default untuk skalar adalah format(12), tetapi dapat diubah.
\end{eulercomment}
\begin{eulerprompt}
>setscalarformat(5); pi
\end{eulerprompt}
\begin{euleroutput}
  3.1416
\end{euleroutput}
\begin{eulerprompt}
>longestformat; pi
\end{eulerprompt}
\begin{euleroutput}
  3.141592653589793
\end{euleroutput}
\begin{eulercomment}
Daftar format keluaran yang penting:

\end{eulercomment}
\begin{eulerttcomment}
 shortestformat shortformat longformat, longestformat
 format(length,digits) goodformat(length)
 fracformat(length)
 defformat
\end{eulerttcomment}
\begin{eulerprompt}
>longestformat; pi,
\end{eulerprompt}
\begin{euleroutput}
  3.141592653589793
\end{euleroutput}
\begin{eulerprompt}
>format(10,5); pi
\end{eulerprompt}
\begin{euleroutput}
    3.14159 
\end{euleroutput}
\begin{eulerprompt}
>defformat; // kembali ke default
\end{eulerprompt}
\begin{eulercomment}
Ada operator singkat yang hanya mencetak satu nilai. Operator "longest"
mencetak semua digit valid suatu bilangan.
\end{eulercomment}
\begin{eulerprompt}
>longest pi^2/2
\end{eulerprompt}
\begin{euleroutput}
        4.934802200544679 
\end{euleroutput}
\begin{eulercomment}
Ada juga operator singkat untuk keluaran pecahan.
\end{eulercomment}
\begin{eulerprompt}
>fraction 1+1/2+1/3+1/4
\end{eulerprompt}
\begin{euleroutput}
  25/12
\end{euleroutput}
\begin{eulercomment}
Karena format internal memakai biner, nilai 0.1 tidak tersimpan eksak.
Galatnya menumpuk sedikit pada perhitungan berikut.
\end{eulercomment}
\begin{eulerprompt}
>longest 0.1+0.1+0.1+0.1+0.1+0.1+0.1+0.1+0.1+0.1-1
\end{eulerprompt}
\begin{euleroutput}
   -1.110223024625157e-16 
\end{euleroutput}
\begin{eulercomment}
Dengan format default Anda tidak akan melihatnya; bilangan yang sangat
kecil ditampilkan sebagai 0.
\end{eulercomment}
\begin{eulerprompt}
>0.1+0.1+0.1+0.1+0.1+0.1+0.1+0.1+0.1+0.1-1
\end{eulerprompt}
\begin{euleroutput}
  0
\end{euleroutput}
\begin{eulercomment}
*** Latihan 7.1
Tampilkan 1/7 dalam 5 digit, dalam 16 digit, lalu sebagai pecahan.
\end{eulercomment}
\begin{eulerprompt}
>format(12,5); 1/7
\end{eulerprompt}
\begin{euleroutput}
      0.14286 
\end{euleroutput}
\begin{eulerprompt}
>longest 1/7
\end{eulerprompt}
\begin{euleroutput}
       0.1428571428571428 
\end{euleroutput}
\begin{eulerprompt}
>defformat; fraction 1/7
\end{eulerprompt}
\begin{euleroutput}
  1/7
\end{euleroutput}
\begin{eulercomment}
*** Latihan 7.2
Hitung 1 + 1/2 + 1/3 + 1/4 + 1/5 + 1/6 sebagai pecahan.
Hasil yang diharapkan: 49/20.
\end{eulercomment}
\begin{eulerprompt}
>fraction sum(1/(1:6))
\end{eulerprompt}
\begin{euleroutput}
  49/20
\end{euleroutput}
\begin{eulercomment}
Penjelasan: format hanya mengatur tampilan, bukan ketelitian internal.
fraction mencari pecahan dengan penyebut kecil yang paling dekat dengan
hasil desimal, sehingga hanya tepat bila hasil aslinya memang rasional
berpenyebut kecil.

(Bersambung ke Bagian 2: Ekspresi, Matematika Simbolik, Fungsi,
Menyelesaikan Persamaan dan Pertidaksamaan.)
* EMT untuk Perhitungan Aljabar (Bagian 2)

Bagian ini mencakup: Ekspresi, Matematika Simbolik, Fungsi, Menyelesaikan
Persamaan, dan Menyelesaikan Pertidaksamaan. Setiap bagian berisi uraian
materi, contoh EMT, soal latihan, dan penjelasan hasil.

Catatan: keluaran perintah latihan baru tidak dicantumkan karena akan muncul
saat Anda meng-ENTER barisnya. Nilai yang diharapkan ditulis di komentar.

\begin{eulercomment}
\eulerheading{Ekspresi}
\begin{eulercomment}
Uraian materi: ekspresi aljabar adalah rangkaian bilangan, variabel, dan
operasi. Nilainya dihitung dengan mensubstitusi nilai variabel.

String atau nama dapat dipakai untuk menyimpan ekspresi matematika yang dapat
dievaluasi EMT. Untuk itu, pakai tanda kurung setelah nama ekspresi. Jika
sebuah string dipakai sebagai ekspresi, gunakan kebiasaan menamainya "fx"
atau "fxy" dan seterusnya. Ekspresi diprioritaskan di atas fungsi.

Variabel global dapat dipakai dalam evaluasi.
\end{eulercomment}
\begin{eulerprompt}
>r:=2; fx:="pi*r^2"; longest fx()
\end{eulerprompt}
\begin{euleroutput}
        12.56637061435917 
\end{euleroutput}
\begin{eulercomment}
Parameter diberikan ke x, y, dan z secara berurutan. Parameter tambahan dapat
diberikan dengan parameter bernama.
\end{eulercomment}
\begin{eulerprompt}
>fx:="a*sin(x)^2"; fx(5,a=-1)
\end{eulerprompt}
\begin{euleroutput}
  -0.919535764538
\end{euleroutput}
\begin{eulercomment}
Perhatikan: ekspresi selalu memakai variabel global, walaupun di dalam
fungsi ada variabel bernama sama. (Jika tidak, evaluasi ekspresi di dalam
fungsi akan membingungkan pengguna fungsi tersebut.)
\end{eulercomment}
\begin{eulerprompt}
>at:=4; function f(expr,x,at) := expr(x); ...
>f("at*x^2",3,5) // menghitung 4*3^2, bukan 5*3^2
\end{eulerprompt}
\begin{euleroutput}
  36
\end{euleroutput}
\begin{eulercomment}
Jika ingin memakai nilai "at" selain nilai global, tambahkan "at=nilai".
\end{eulercomment}
\begin{eulerprompt}
>at:=4; function f(expr,x,a) := expr(x,at=a); ...
>f("at*x^2",3,5)
\end{eulerprompt}
\begin{euleroutput}
  45
\end{euleroutput}
\begin{eulercomment}
Koleksi pemanggilan (call collection) dapat memuat ekspresi. Contoh di atas
dapat ditulis sebagai berikut.
\end{eulercomment}
\begin{eulerprompt}
>at:=4; function f(expr,x) := expr(x); ...
>f(\{\{"at*x^2",at=5\}\},3)
\end{eulerprompt}
\begin{euleroutput}
  45
\end{euleroutput}
\begin{eulercomment}
Ekspresi dalam x sering dipakai seperti fungsi. Perhatikan bahwa mendefinisikan
fungsi dengan nama yang sama dengan ekspresi simbolik global akan menghapus
variabel tersebut, untuk menghindari kerancuan.
\end{eulercomment}
\begin{eulerprompt}
>f &= 5*x;
>function f(x) := 6*x;
>f(2)
\end{eulerprompt}
\begin{euleroutput}
  12
\end{euleroutput}
\begin{eulercomment}
Menurut kebiasaan, ekspresi simbolik atau numerik dinamai fx, fxy, dst.
Pola nama ini sebaiknya tidak dipakai untuk fungsi.
\end{eulercomment}
\begin{eulerprompt}
>fx &= diff(x^x,x); $&fx
\end{eulerprompt}
\begin{eulercomment}
Bentuk khusus ekspresi memungkinkan variabel apa pun sebagai parameter tanpa
nama, bukan hanya "x", "y". Awali ekspresi dengan "@(variabel) ...".
\end{eulercomment}
\begin{eulerprompt}
>"@(a,b) a^2+b^2", %(4,5)
\end{eulerprompt}
\begin{euleroutput}
  @(a,b) a^2+b^2
  41
\end{euleroutput}
\begin{eulercomment}
Cara paling dasar mendefinisikan fungsi sederhana adalah menyimpan rumusnya
dalam ekspresi. Jika variabel utamanya x, ekspresi dapat dievaluasi seperti
fungsi. Variabel global terlihat saat evaluasi.
\end{eulercomment}
\begin{eulerprompt}
>fx &= x^3-a*x;  ...
>a=1.2; fx(0.5)
\end{eulerprompt}
\begin{euleroutput}
  -0.475
\end{euleroutput}
\begin{eulercomment}
Variabel lain dalam ekspresi dapat ditentukan saat evaluasi dengan parameter
bernama.
\end{eulercomment}
\begin{eulerprompt}
>fx(0.5,a=1.1)
\end{eulerprompt}
\begin{euleroutput}
  -0.425
\end{euleroutput}
\begin{eulercomment}
Ekspresi tidak harus simbolik. Ini diperlukan jika ekspresi memuat fungsi
yang hanya dikenal kernel numerik, bukan Maxima.

*** Latihan 8.1
Definisikan fx:="a*x\textasciicircum{}2+b" lalu hitung untuk x=2, a=3, b=1.
Hasil yang diharapkan: 3*4+1 = 13.
\end{eulercomment}
\begin{eulerprompt}
>fx:="a*x^2+b"; fx(2,a=3,b=1)
\end{eulerprompt}
\begin{euleroutput}
  13
\end{euleroutput}
\begin{eulercomment}
*** Latihan 8.2
Evaluasi "@(a,b) a\textasciicircum{}2+b\textasciicircum{}2" pada (3,4). Hasil yang diharapkan: 25.
\end{eulercomment}
\begin{eulerprompt}
>"@(a,b) a^2+b^2", %(3,4)
\end{eulerprompt}
\begin{euleroutput}
  @(a,b) a^2+b^2
  25
\end{euleroutput}
\begin{eulercomment}
*** Latihan 8.3
Definisikan fx \&= x\textasciicircum{}3-3*x lalu hitung fx(2) dan fx(-1). Hasil yang
diharapkan: 2 dan 2.
\end{eulercomment}
\begin{eulerprompt}
>fx &= x^3-3*x; fx(2), fx(-1)
\end{eulerprompt}
\begin{euleroutput}
  2
  2
\end{euleroutput}
\begin{eulercomment}
Penjelasan: ekspresi dievaluasi dengan nilai parameter yang diberikan, dan
variabel yang tidak diberikan diambil dari variabel global.

\begin{eulercomment}
\eulerheading{Matematika Simbolik}
\begin{eulercomment}
Uraian materi: manipulasi simbolik bekerja pada rumus eksak, tanpa
pembulatan. Operasi pokoknya: menjabarkan (expand), memfaktorkan (factor),
menyederhanakan, mendiferensialkan, dan mengintegralkan.

EMT melakukan matematika simbolik dengan bantuan Maxima. Matematika simbolik
terintegrasi dengan Euler melalui \&. Ekspresi yang diawali \& adalah ekspresi
simbolik, dievaluasi dan dicetak oleh Maxima.

Pertama, Maxima memiliki aritmetika "tak hingga" yang dapat menangani
bilangan sangat besar.
\end{eulercomment}
\begin{eulerprompt}
>$&44!
\end{eulerprompt}
\begin{eulercomment}
Dengan demikian, hasil besar dapat dihitung secara eksak. Hitung

\end{eulercomment}
\begin{eulerformula}
\[
C(44,10) = \frac{44!}{34! \cdot 10!}
\]
\end{eulerformula}
\begin{eulerprompt}
>$& 44!/(34!*10!) // nilai C(44,10)
\end{eulerprompt}
\begin{eulercomment}
Tentu Maxima punya fungsi yang lebih efisien untuk ini.
\end{eulercomment}
\begin{eulerprompt}
>$binomial(44,10) //menghitung C(44,10) menggunakan fungsi binomial()
\end{eulerprompt}
\begin{eulercomment}
Untuk mempelajari fungsi tertentu, klik ganda pada fungsi tersebut. Berikut
juga berlaku:

\end{eulercomment}
\begin{eulerformula}
\[
C(x,3)=\frac{x!}{(x-3)!3!}=\frac{(x-2)(x-1)x}{6}
\]
\end{eulerformula}
\begin{eulerprompt}
>$binomial(x,3) // C(x,3)
\end{eulerprompt}
\begin{eulercomment}
Untuk mengganti x dengan nilai tertentu gunakan "with".
\end{eulercomment}
\begin{eulerprompt}
>$&binomial(x,3) with x=10 // substitusi x=10 ke C(x,3)
\end{eulerprompt}
\begin{eulercomment}
Ekspresi simbolik dicetak Maxima dalam bentuk 2D. Jika LaTeX terpasang,
ekspresi simbolik dapat dicetak dengan LaTeX: awali dengan \textdollar{} di depan \&
(atau \& boleh dihilangkan). Jangan jalankan perintah Maxima dengan \textdollar{} jika
LaTeX tidak terpasang.
\end{eulercomment}
\begin{eulerprompt}
>$(3+x)/(x^2+1)
\end{eulerprompt}
\begin{eulercomment}
Ekspresi simbolik di-parse oleh Euler. Jika perlu sintaks kompleks dalam
satu ekspresi, bungkus ekspresi dengan "...". Menggunakan lebih dari satu
ekspresi sederhana dimungkinkan tetapi sangat tidak dianjurkan.
\end{eulercomment}
\begin{eulerprompt}
>&"v := 5; v^2"
\end{eulerprompt}
\begin{euleroutput}
  
                                    25
  
\end{euleroutput}
\begin{eulercomment}
Ekspresi simbolik dapat dipakai dalam program, tetapi harus dibungkus tanda
kutip.
\end{eulercomment}
\begin{eulerprompt}
>$&expand((1+x)^4), $&factor(diff(%,x)) // diff: turunan, factor: faktor
\end{eulerprompt}
\begin{eulercomment}
Lagi, \% merujuk ke hasil sebelumnya.

Agar mudah, simpan solusi ke variabel simbolik. Variabel simbolik
didefinisikan dengan "\&=".
\end{eulercomment}
\begin{eulerprompt}
>fx &= (x+1)/(x^4+1); $&fx
\end{eulerprompt}
\begin{eulercomment}
Ekspresi simbolik dapat dipakai dalam ekspresi simbolik lain.
\end{eulercomment}
\begin{eulerprompt}
>$&factor(diff(fx,x))
\end{eulerprompt}
\begin{eulercomment}
Perintah Maxima dapat dimasukkan langsung dengan awalan "::". Sintaks Maxima
disesuaikan dengan sintaks EMT ("mode kompatibilitas").
\end{eulercomment}
\begin{eulerprompt}
>&factor(20!)
\end{eulerprompt}
\begin{euleroutput}
  
                           2432902008176640000
  
\end{euleroutput}
\begin{eulerprompt}
>::: factor(10!)
\end{eulerprompt}
\begin{euleroutput}
  
                                 8  4  2
                                2  3  5  7
  
\end{euleroutput}
\begin{eulerprompt}
>:: factor(20!)
\end{eulerprompt}
\begin{euleroutput}
  
                          18  8  4  2
                         2   3  5  7  11 13 17 19
  
\end{euleroutput}
\begin{eulercomment}
Jika Anda ahli Maxima, Anda dapat memakai sintaks asli Maxima dengan ":::".
\end{eulercomment}
\begin{eulerprompt}
>::: av:g$ av^2;
\end{eulerprompt}
\begin{euleroutput}
  
                                     2
                                    g
  
\end{euleroutput}
\begin{eulerprompt}
>fx &= x^3*exp(x), $fx
\end{eulerprompt}
\begin{euleroutput}
  
                                   3  x
                                  x  E
  
\end{euleroutput}
\begin{eulercomment}
Variabel seperti ini dapat dipakai dalam ekspresi simbolik lain. Perhatikan
bahwa pada perintah berikut ruas kanan \&= dievaluasi sebelum penugasan.
\end{eulercomment}
\begin{eulerprompt}
>&(fx with x=5), $%, &float(%)
\end{eulerprompt}
\begin{euleroutput}
  
                                       5
                                  125 E
  
  
                            18551.64488782208
  
\end{euleroutput}
\begin{eulerprompt}
>fx(5)
\end{eulerprompt}
\begin{euleroutput}
  18551.6448878
\end{euleroutput}
\begin{eulercomment}
Untuk mengevaluasi ekspresi pada nilai variabel tertentu, gunakan operator
"with". Perintah berikut juga menunjukkan bahwa Maxima dapat mengevaluasi
ekspresi secara numerik dengan float().
\end{eulercomment}
\begin{eulerprompt}
>&(fx with x=10)-(fx with x=5), &float(%)
\end{eulerprompt}
\begin{euleroutput}
  
                                  10        5
                            1000 E   - 125 E
  
  
                           2.20079141499189e+7
  
\end{euleroutput}
\begin{eulerprompt}
>$factor(diff(fx,x,2))
\end{eulerprompt}
\begin{eulercomment}
Untuk mendapat kode LaTeX suatu ekspresi, gunakan perintah tex.
\end{eulercomment}
\begin{eulerprompt}
>tex(fx)
\end{eulerprompt}
\begin{euleroutput}
  x^3\(\backslash\),e^\{x\}
\end{euleroutput}
\begin{eulercomment}
Ekspresi simbolik dapat dievaluasi seperti ekspresi numerik.
\end{eulercomment}
\begin{eulerprompt}
>fx(0.5)
\end{eulerprompt}
\begin{euleroutput}
  0.206090158838
\end{euleroutput}
\begin{eulercomment}
Pada ekspresi simbolik, ini tidak berlaku karena Maxima tidak
mendukungnya. Gunakan sintaks "with" (bentuk lebih baik dari perintah
at(...) Maxima).
\end{eulercomment}
\begin{eulerprompt}
>$&fx with x=1/2
\end{eulerprompt}
\begin{eulercomment}
Penugasannya juga dapat berupa simbolik.
\end{eulercomment}
\begin{eulerprompt}
>$&fx with x=1+t
\end{eulerprompt}
\begin{eulercomment}
Perintah solve menyelesaikan ekspresi simbolik untuk suatu variabel di
Maxima. Hasilnya berupa vektor solusi.
\end{eulercomment}
\begin{eulerprompt}
>$&solve(x^2+x=4,x)
\end{eulerprompt}
\begin{eulercomment}
Bandingkan dengan perintah numerik "solve" di Euler, yang memerlukan nilai
awal dan, opsional, nilai target.
\end{eulercomment}
\begin{eulerprompt}
>solve("x^2+x",1,y=4)
\end{eulerprompt}
\begin{euleroutput}
  1.56155281281
\end{euleroutput}
\begin{eulercomment}
Nilai numerik solusi simbolik dapat dihitung dengan mengevaluasi hasil
simbolik. Euler akan melewati penugasan x= dst.
\end{eulercomment}
\begin{eulerprompt}
>sol &= solve(x^2+2*x=4,x); $&sol, sol(), $&float(sol)
\end{eulerprompt}
\begin{euleroutput}
  [-3.23607,  1.23607]
\end{euleroutput}
\begin{eulercomment}
Untuk mengambil solusi simbolik tertentu, gunakan "with" dan indeks.
\end{eulercomment}
\begin{eulerprompt}
>$&solve(x^2+x=1,x), x2 &= x with %[2]; $&x2
\end{eulerprompt}
\begin{eulercomment}
Untuk menyelesaikan sistem persamaan, gunakan vektor persamaan. Hasilnya
berupa vektor solusi.
\end{eulercomment}
\begin{eulerprompt}
>sol &= solve([x+y=3,x^2+y^2=5],[x,y]); $&sol, $&x*y with sol[1]
\end{eulerprompt}
\begin{eulercomment}
Ekspresi simbolik dapat memiliki flag yang menunjukkan perlakuan khusus di
Maxima. Flag ditambahkan dengan "\textbar{}" (bentuk lebih baik dari "ev(...,flags)").
\end{eulercomment}
\begin{eulerprompt}
>$& diff((x^3-1)/(x+1),x) //turunan bentuk pecahan
>$& diff((x^3-1)/(x+1),x) | ratsimp //menyederhanakan pecahan
>$&factor(%)
\end{eulerprompt}
\begin{eulercomment}
*** Latihan 9.1
Hitung C(10,3) dengan dua cara. Hasil yang diharapkan: 120.
\end{eulercomment}
\begin{eulerprompt}
>$&binomial(10,3)
>$&10!/(7!*3!)
\end{eulerprompt}
\begin{eulercomment}
*** Latihan 9.2
Jabarkan (x+2)\textasciicircum{}5.
Hasil yang diharapkan: x\textasciicircum{}5+10x\textasciicircum{}4+40x\textasciicircum{}3+80x\textasciicircum{}2+80x+32.
\end{eulercomment}
\begin{eulerprompt}
>$&expand((x+2)^5)
\end{eulerprompt}
\begin{eulercomment}
*** Latihan 9.3
Tentukan turunan x\textasciicircum{}3*e\textasciicircum{}x dan faktorkan.
Hasil yang diharapkan: x\textasciicircum{}2 e\textasciicircum{}x (x+3).
\end{eulercomment}
\begin{eulerprompt}
>$&factor(diff(x^3*exp(x),x))
\end{eulerprompt}
\begin{eulercomment}
*** Latihan 9.4
Hitung 30! secara eksak dan faktorisasi prima 100.
100 = 2\textasciicircum{}2 * 5\textasciicircum{}2.
\end{eulercomment}
\begin{eulerprompt}
>$&30!
>$&factor(100)
\end{eulerprompt}
\begin{eulercomment}
*** Latihan 9.5
Sederhanakan (x\textasciicircum{}2-1)/(x-1). Hasil yang diharapkan: x+1.
\end{eulercomment}
\begin{eulerprompt}
>$&ratsimp((x^2-1)/(x-1))
\end{eulerprompt}
\begin{eulercomment}
Penjelasan: Maxima menghitung secara eksak (bilangan bulat besar, akar, dan
pecahan tidak dibulatkan). \&float() mengubah hasil eksak menjadi desimal.

\begin{eulercomment}
\eulerheading{Fungsi}
\begin{eulercomment}
Uraian materi: fungsi f memetakan setiap x pada tepat satu nilai f(x).
Di EMT, fungsi adalah program yang didefinisikan dengan perintah "function",
berupa fungsi satu baris atau multi-baris.

Fungsi satu baris dapat berupa numerik atau simbolik. Fungsi satu baris
numerik didefinisikan dengan ":=".
\end{eulercomment}
\begin{eulerprompt}
>function f(x) := x*sqrt(x^2+1)
\end{eulerprompt}
\begin{eulercomment}
Fungsi dapat dievaluasi seperti fungsi bawaan Euler.
\end{eulercomment}
\begin{eulerprompt}
>f(2)
\end{eulerprompt}
\begin{euleroutput}
  4.472135955
\end{euleroutput}
\begin{eulercomment}
Fungsi ini juga bekerja untuk vektor, mengikuti bahasa matriks Euler,
karena ekspresi di dalamnya sudah divektorkan.
\end{eulercomment}
\begin{eulerprompt}
>f(0:0.1:1)
\end{eulerprompt}
\begin{euleroutput}
  [0,  0.100499,  0.203961,  0.313209,  0.430813,  0.559017,  0.699714,
  0.854459,  1.0245,  1.21083,  1.41421]
\end{euleroutput}
\begin{eulercomment}
Berbeda dengan ekspresi, nama fungsi harus diberikan sebagai string pada
perintah seperti solve.
\end{eulercomment}
\begin{eulerprompt}
>solve("f",1,y=1)
\end{eulerprompt}
\begin{euleroutput}
  0.786151377757
\end{euleroutput}
\begin{eulercomment}
Untuk menimpa fungsi bawaan, tambahkan kata kunci "overwrite". Menimpa
fungsi bawaan berbahaya dan dapat mengganggu fungsi lain. Fungsi bawaan di
inti Euler masih dapat dipanggil dengan "\_...".
\end{eulercomment}
\begin{eulerprompt}
>function overwrite sin (x) := _sin(x°) // mendefinisikan ulang sin dalam derajat
>sin(45)
\end{eulerprompt}
\begin{euleroutput}
  Variable x not found!
  Use global or local variables defined in function sin.
  Error in degree
  Try "trace errors" to inspect local variables after errors.
  sin:
      useglobal; return _sin(x°)  
  Error in:
  sin(45) ...
         ^
\end{euleroutput}
\begin{eulercomment}
Sebaiknya hapus redefinisi ini.
\end{eulercomment}
\begin{eulerprompt}
>forget sin; sin(pi/4)
\end{eulerprompt}
\begin{euleroutput}
  0.707106781187
\end{euleroutput}
\eulersubheading{Parameter Default}
\begin{eulercomment}
Fungsi numerik dapat memiliki parameter default.
\end{eulercomment}
\begin{eulerprompt}
>function f(x,a=1) := a*x^2
\end{eulerprompt}
\begin{eulercomment}
Jika parameter dihilangkan, nilai default dipakai.
\end{eulercomment}
\begin{eulerprompt}
>f(4)
\end{eulerprompt}
\begin{euleroutput}
  16
\end{euleroutput}
\begin{eulercomment}
Jika diberikan, nilai default ditimpa.
\end{eulercomment}
\begin{eulerprompt}
>f(4,5)
\end{eulerprompt}
\begin{euleroutput}
  80
\end{euleroutput}
\begin{eulercomment}
Parameter bernama juga menimpanya. Ini dipakai banyak fungsi Euler seperti
plot2d dan plot3d.
\end{eulercomment}
\begin{eulerprompt}
>f(4,a=1)
\end{eulerprompt}
\begin{euleroutput}
  16
\end{euleroutput}
\begin{eulercomment}
Jika variabel bukan parameter, variabel itu harus global. Fungsi satu baris
dapat melihat variabel global.
\end{eulercomment}
\begin{eulerprompt}
>function f(x) := a*x^2
>a=6; f(2)
\end{eulerprompt}
\begin{euleroutput}
  24
\end{euleroutput}
\begin{eulercomment}
Parameter bernama menimpa nilai global. Jika argumen tidak ada dalam daftar
parameter, harus dideklarasikan dengan ":=".
\end{eulercomment}
\begin{eulerprompt}
>f(2,a:=5)
\end{eulerprompt}
\begin{euleroutput}
  20
\end{euleroutput}
\begin{eulercomment}
Fungsi simbolik didefinisikan dengan "\&=". Fungsi ini didefinisikan di Euler
dan Maxima, dan bekerja di kedua dunia. Ekspresi pendefinisinya diproses
Maxima sebelum didefinisikan.
\end{eulercomment}
\begin{eulerprompt}
>function g(x) &= x^3-x*exp(-x); $&g(x)
\end{eulerprompt}
\begin{eulercomment}
Fungsi simbolik dapat dipakai dalam ekspresi simbolik.
\end{eulercomment}
\begin{eulerprompt}
>$&diff(g(x),x), $&% with x=4/3
\end{eulerprompt}
\begin{eulercomment}
Fungsi simbolik juga dapat dipakai dalam ekspresi numerik, selama EMT dapat
menafsirkan seluruh isi fungsi.
\end{eulercomment}
\begin{eulerprompt}
>g(5+g(1))
\end{eulerprompt}
\begin{euleroutput}
  178.635099908
\end{euleroutput}
\begin{eulercomment}
Fungsi simbolik dapat dipakai untuk mendefinisikan fungsi atau ekspresi
simbolik lain.
\end{eulercomment}
\begin{eulerprompt}
>function G(x) &= factor(integrate(g(x),x)); $&G(c) // integrate: mengintegralkan
>solve(&g(x),0.5)
\end{eulerprompt}
\begin{euleroutput}
  0.703467422498
\end{euleroutput}
\begin{eulercomment}
Berikut juga bekerja, karena Euler memakai ekspresi simbolik dalam fungsi g
jika tidak menemukan variabel simbolik g dan ada fungsi simbolik g.
\end{eulercomment}
\begin{eulerprompt}
>solve(&g,0.5)
\end{eulerprompt}
\begin{euleroutput}
  0.703467422498
\end{euleroutput}
\begin{eulerprompt}
>function P(x,n) &= (2*x-1)^n; $&P(x,n)
>function Q(x,n) &= (x+2)^n; $&Q(x,n)
>$&P(x,4), $&expand(%)
>P(3,4)
\end{eulerprompt}
\begin{euleroutput}
  625
\end{euleroutput}
\begin{eulerprompt}
>$&P(x,4)+ Q(x,3), $&expand(%)
>$&P(x,4)-Q(x,3), $&expand(%), $&factor(%)
>$&P(x,4)*Q(x,3), $&expand(%), $&factor(%)
>$&P(x,4)/Q(x,1), $&expand(%), $&factor(%)
>function f(x) &= x^3-x; $&f(x)
\end{eulerprompt}
\begin{eulercomment}
Dengan \&= fungsi bersifat simbolik dan dapat dipakai dalam ekspresi simbolik
lain.
\end{eulercomment}
\begin{eulerprompt}
>$&integrate(f(x),x)
\end{eulerprompt}
\begin{eulercomment}
Dengan := fungsi bersifat numerik. Contoh yang baik adalah integral tentu

\end{eulercomment}
\begin{eulerformula}
\[
f(x) = \int_1^x t^t \, dt,
\]
\end{eulerformula}
\begin{eulercomment}
yang tidak dapat dievaluasi secara simbolik.

Jika fungsi didefinisikan ulang dengan kata kunci "map", fungsi itu dapat
dipakai untuk vektor x. Secara internal, fungsi dipanggil untuk setiap nilai
x dan hasilnya disimpan dalam vektor.
\end{eulercomment}
\begin{eulerprompt}
>function map f(x) := integrate("x^x",1,x)
>f(0:0.5:2)
\end{eulerprompt}
\begin{euleroutput}
  [-0.783431,  -0.410816,  0,  0.676863,  2.05045]
\end{euleroutput}
\begin{eulercomment}
Fungsi dapat memiliki nilai default untuk parameter.
\end{eulercomment}
\begin{eulerprompt}
>function mylog (x,base=10) := ln(x)/ln(base);
\end{eulerprompt}
\begin{eulercomment}
Sekarang fungsi dapat dipanggil dengan atau tanpa parameter "base".
\end{eulercomment}
\begin{eulerprompt}
>mylog(100), mylog(2^6.7,2)
\end{eulerprompt}
\begin{euleroutput}
  2
  6.7
\end{euleroutput}
\begin{eulerprompt}
>mylog(E^2,base=E)
\end{eulerprompt}
\begin{euleroutput}
  2
\end{euleroutput}
\begin{eulercomment}
Sering kita ingin memakai fungsi untuk vektor di satu tempat dan untuk
elemen tunggal di tempat lain. Ini mungkin dengan parameter vektor.
\end{eulercomment}
\begin{eulerprompt}
>function f([a,b]) &= a^2+b^2-a*b+b; $&f(a,b), $&f(x,y)
\end{eulerprompt}
\begin{eulercomment}
Fungsi simbolik seperti ini dapat dipakai untuk variabel simbolik, dan juga
untuk vektor numerik.
\end{eulercomment}
\begin{eulerprompt}
>v=[3,4]; f(v)
\end{eulerprompt}
\begin{euleroutput}
  17
\end{euleroutput}
\begin{eulercomment}
Ada juga fungsi simbolik murni yang tidak dapat dipakai secara numerik.
\end{eulercomment}
\begin{eulerprompt}
>function lapl(expr,x,y) &&= diff(expr,x,2)+diff(expr,y,2)//turunan parsial kedua
\end{eulerprompt}
\begin{euleroutput}
  
                   diff(expr, y, 2) + diff(expr, x, 2)
  
\end{euleroutput}
\begin{eulerprompt}
>$&realpart((x+I*y)^4), $&lapl(%,x,y)
>function f(x,y) &= factor(lapl((x+y^2)^5,x,y)); $&f(x,y)
\end{eulerprompt}
\begin{eulercomment}
Ringkasan:

- \&= mendefinisikan fungsi simbolik,
- := mendefinisikan fungsi numerik,
- \&\&= mendefinisikan fungsi simbolik murni.

*** Latihan 10.1
Definisikan f(x) := x\textasciicircum{}2-3x+2. Hitung f(5) dan f(0:5).
Hasil yang diharapkan: 12 dan [2,0,0,2,6,12].
\end{eulercomment}
\begin{eulerprompt}
>function f(x) := x^2-3*x+2
>f(5), f(0:5)
\end{eulerprompt}
\begin{euleroutput}
  12
  [2,  0,  0,  2,  6,  12]
\end{euleroutput}
\begin{eulercomment}
*** Latihan 10.2
Definisikan fungsi simbolik h(x) \&= x\textasciicircum{}3-3x. Tentukan turunan dan
integralnya. Hasil yang diharapkan: 3x\textasciicircum{}2-3 dan x\textasciicircum{}4/4 - 3x\textasciicircum{}2/2.
\end{eulercomment}
\begin{eulerprompt}
>function h(x) &= x^3-3*x; $&h(x)
>$&diff(h(x),x), $&integrate(h(x),x)
\end{eulerprompt}
\begin{eulercomment}
*** Latihan 10.3
Dengan mylog, hitung log basis 10 dari 1000 dan log basis 2 dari 8.
Hasil yang diharapkan: 3 dan 3.
\end{eulercomment}
\begin{eulerprompt}
>mylog(1000), mylog(8,2)
\end{eulerprompt}
\begin{euleroutput}
  3
  3
\end{euleroutput}
\begin{eulercomment}
*** Latihan 10.4
Dengan P dan Q di atas, jabarkan P(x,2)*Q(x,1).
Hasil yang diharapkan: 4x\textasciicircum{}3+4x\textasciicircum{}2-7x+2.
\end{eulercomment}
\begin{eulerprompt}
>$&expand(P(x,2)*Q(x,1))
\end{eulerprompt}
\begin{eulercomment}
*** Latihan 10.5
Tentukan nilai maksimum dan minimum lokal h(x)=x\textasciicircum{}3-3x dengan mencari titik
stasioner: h'(x)=0 memberi x=-1 dan x=1.
\end{eulercomment}
\begin{eulerprompt}
>$&solve(diff(h(x),x)=0,x)
>h(-1), h(1)  // seharusnya 2 dan -2
\end{eulerprompt}
\begin{euleroutput}
  2
  -2
\end{euleroutput}
\begin{eulercomment}
Penjelasan: "\&=" membuat fungsi yang dapat diturunkan/diintegralkan secara
eksak, sedangkan ":=" membuat fungsi numerik murni yang cukup untuk
pengevaluasian nilai.

\begin{eulercomment}
\eulerheading{Menyelesaikan Persamaan}
\begin{eulercomment}
Uraian materi: akar persamaan f(x)=0 adalah nilai x yang membuat f(x)=0.
Untuk persamaan kuadrat ax\textasciicircum{}2+bx+c=0 berlaku rumus

\end{eulercomment}
\begin{eulerformula}
\[
x_{1,2}=\frac{-b\pm\sqrt{b^2-4ac}}{2a}
\]
\end{eulerformula}
\begin{eulercomment}
Persamaan dapat diselesaikan secara numerik dan simbolik.

Untuk persamaan satu variabel dipakai fungsi solve(), yang memerlukan nilai
awal. Secara internal solve() memakai metode secant.
\end{eulercomment}
\begin{eulerprompt}
>solve("x^2-2",1)
\end{eulerprompt}
\begin{euleroutput}
  1.41421356237
\end{euleroutput}
\begin{eulercomment}
Ini juga bekerja untuk ekspresi simbolik.
\end{eulercomment}
\begin{eulerprompt}
>$&solve(x^2=2,x)
>$&solve(x^2-2,x)
>$&solve(a*x^2+b*x+c=0,x)
>$&solve([a*x+b*y=c,d*x+e*y=f],[x,y])
>px &= 4*x^8+x^7-x^4-x; $&px
\end{eulerprompt}
\begin{eulercomment}
Sekarang cari titik tempat polinomial bernilai 2. Pada solve(), nilai target
default y=0 dapat diubah dengan variabel bernama. Gunakan y=2 dan periksa
dengan mengevaluasi polinomial pada hasil sebelumnya.
\end{eulercomment}
\begin{eulerprompt}
>solve(px,1,y=2), px(%)
\end{eulerprompt}
\begin{euleroutput}
  0.966715594851
  2
\end{euleroutput}
\begin{eulercomment}
Menyelesaikan ekspresi simbolik secara simbolik menghasilkan daftar solusi.
\end{eulercomment}
\begin{eulerprompt}
>sol &= solve(x^2-x-1,x); $&sol
\end{eulerprompt}
\begin{eulercomment}
Cara termudah mendapat nilai numerik adalah mengevaluasi solusi seperti
ekspresi.
\end{eulercomment}
\begin{eulerprompt}
>longest sol()
\end{eulerprompt}
\begin{euleroutput}
      -0.6180339887498949       1.618033988749895 
\end{euleroutput}
\begin{eulercomment}
Untuk memakai solusi secara simbolik dalam ekspresi lain, cara termudah
adalah "with".
\end{eulercomment}
\begin{eulerprompt}
>$&x^2 with sol[1], $&expand(x^2-x-1 with sol[2])
\end{eulerprompt}
\begin{eulercomment}
Sistem persamaan secara simbolik: gunakan vektor persamaan. Jawabannya
berupa daftar dari daftar persamaan.
\end{eulercomment}
\begin{eulerprompt}
>$&solve([x+y=2,x^3+2*y+x=4],[x,y])
\end{eulerprompt}
\begin{eulercomment}
Fungsi f() dapat melihat variabel global, tetapi sering kita ingin
parameter lokal.

\end{eulercomment}
\begin{eulerformula}
\[
a^x-x^a = 0.1
\]
\end{eulerformula}
\begin{eulercomment}
dengan a=3.
\end{eulercomment}
\begin{eulerprompt}
>function f(x,a) := x^a-a^x;
\end{eulerprompt}
\begin{eulercomment}
Satu cara meneruskan parameter tambahan ke f() adalah memakai daftar berisi
nama fungsi dan parameternya (cara lain adalah parameter titik koma).
\end{eulercomment}
\begin{eulerprompt}
>solve(\{\{"f",3\}\},2,y=0.1)
\end{eulerprompt}
\begin{euleroutput}
  2.54116291558
\end{euleroutput}
\begin{eulercomment}
Ini juga bekerja untuk ekspresi, tetapi harus memakai elemen daftar bernama.
\end{eulercomment}
\begin{eulerprompt}
>solve(\{\{"x^a-a^x",a=3\}\},2,y=0.1)
\end{eulerprompt}
\begin{euleroutput}
  2.54116291558
\end{euleroutput}
\begin{eulercomment}
*** Latihan 11.1
Selesaikan x\textasciicircum{}2-7x+10=0 secara simbolik. Hasil yang diharapkan: x=2 dan x=5.
\end{eulercomment}
\begin{eulerprompt}
>$&solve(x^2-7*x+10=0,x)
\end{eulerprompt}
\begin{eulercomment}
*** Latihan 11.2
Selesaikan x\textasciicircum{}3-2x-5=0 secara numerik dengan nilai awal 2.
Hasil yang diharapkan: sekitar 2.09455.
\end{eulercomment}
\begin{eulerprompt}
>solve("x^3-2*x-5",2)
\end{eulerprompt}
\begin{euleroutput}
  2.09455148154
\end{euleroutput}
\begin{eulercomment}
*** Latihan 11.3
Selesaikan sistem 2x+y=7 dan x-y=2. Hasil yang diharapkan: x=3, y=1.
\end{eulercomment}
\begin{eulerprompt}
>$&solve([2*x+y=7,x-y=2],[x,y])
\end{eulerprompt}
\begin{eulercomment}
*** Latihan 11.4
Tentukan x positif yang memenuhi x\textasciicircum{}2=9 memakai solve numerik dengan
target y=9 (nilai awal 1). Hasil yang diharapkan: 3.
\end{eulercomment}
\begin{eulerprompt}
>solve("x^2",1,y=9)
\end{eulerprompt}
\begin{euleroutput}
  3
\end{euleroutput}
\begin{eulercomment}
*** Latihan 11.5
Gunakan rumus kuadrat secara simbolik untuk x\textasciicircum{}2+bx+1=0 dan tentukan nilai b
agar diskriminan nol. Diskriminan b\textasciicircum{}2-4=0 memberi b=2 atau b=-2.
\end{eulercomment}
\begin{eulerprompt}
>$&solve(x^2+b*x+1=0,x)
>$&solve(b^2-4=0,b)
\end{eulerprompt}
\begin{eulercomment}
Penjelasan: solve numerik mencari akar terdekat dari nilai awal (hanya satu
akar per pemanggilan), sedangkan solve simbolik memberi semua akar eksak.

\begin{eulercomment}
\eulerheading{Menyelesaikan Pertidaksamaan}
\begin{eulercomment}
Uraian materi: pertidaksamaan memuat tanda \textless{}, \textgreater{}, \textless{}=, \textgreater{}=. Himpunan
penyelesaiannya berupa interval atau gabungan interval pada garis bilangan.

Untuk menyelesaikan pertidaksamaan, EMT tidak dapat melakukannya sendiri,
melainkan dengan bantuan Maxima, yaitu secara eksak (simbolik). Perintah
Maxima yang dipakai adalah fourier\_elim(), yang harus dipanggil dengan
perintah "load(fourier\_elim)" terlebih dahulu.
\end{eulercomment}
\begin{eulerprompt}
>&load(fourier_elim)
\end{eulerprompt}
\begin{euleroutput}
  
          C:/Program Files/Euler x64/maxima/share/maxima/5.35.1/share/f\(\backslash\)
  ourier_elim/fourier_elim.lisp
  
\end{euleroutput}
\begin{eulerprompt}
>$&fourier_elim([x^2 - 1>0],[x]) // x^2-1 > 0
>$&fourier_elim([x^2 - 1<0],[x]) // x^2-1 < 0
>$&fourier_elim([x^2 - 1 # 0],[x]) // x^2-1 <> 0
>$&fourier_elim([x # 6],[x])
>$&fourier_elim([x < 1, x > 1],[x]) // tidak memiliki penyelesaian
>$&fourier_elim([minf < x, x < inf],[x]) // solusinya R
>$&fourier_elim([x^3 - 1 > 0],[x])
>$&fourier_elim([cos(x) < 1/2],[x]) // gagal (fungsi trigonometri tidak didukung)
>$&fourier_elim([y-x < 5, x - y < 7, 10 < y],[x,y]) // sistem pertidaksamaan
>$&fourier_elim([y-x < 5, x - y < 7, 10 < y],[y,x])
>$&fourier_elim((x + y < 5) and (x - y >8),[x,y])
>$&fourier_elim(((x + y < 5) and x < 1) or  (x - y >8),[x,y])
>&fourier_elim([max(x,y) > 6, x # 8, abs(y-1) > 12],[x,y])
\end{eulerprompt}
\begin{euleroutput}
  
          [6 < x, x < 8, y < - 11] or [8 < x, y < - 11]
   or [x < 8, 13 < y] or [x = y, 13 < y] or [8 < x, x < y, 13 < y]
   or [y < x, 13 < y]
  
\end{euleroutput}
\begin{eulerprompt}
>$&fourier_elim([(x+6)/(x-9) <= 6],[x])
\end{eulerprompt}
\begin{eulercomment}
*** Latihan 12.1
Selesaikan x\textasciicircum{}2-5x+6\textless{}0. Hasil yang diharapkan: 2 \textless{} x \textless{} 3.
\end{eulercomment}
\begin{eulerprompt}
>$&fourier_elim([x^2-5*x+6<0],[x])
\end{eulerprompt}
\begin{eulercomment}
*** Latihan 12.2
Selesaikan x\textasciicircum{}2-4\textgreater{}=0. Hasil yang diharapkan: x\textless{}=-2 atau x\textgreater{}=2.
\end{eulercomment}
\begin{eulerprompt}
>$&fourier_elim([x^2-4>=0],[x])
\end{eulerprompt}
\begin{eulercomment}
*** Latihan 12.3
Selesaikan sistem 1\textless{}x, x\textless{}5, x#3 (x tidak sama dengan 3).
Hasil yang diharapkan: 1\textless{}x\textless{}3 atau 3\textless{}x\textless{}5.
\end{eulercomment}
\begin{eulerprompt}
>$&fourier_elim([1<x,x<5,x#3],[x])
\end{eulerprompt}
\begin{eulercomment}
*** Latihan 12.4
Selesaikan \textbar{}x-2\textbar{}\textless{}3. Hasil yang diharapkan: -1 \textless{} x \textless{} 5.
\end{eulercomment}
\begin{eulerprompt}
>$&fourier_elim([abs(x-2)<3],[x])
\end{eulerprompt}
\begin{eulercomment}
Penjelasan: fourier\_elim() mengeliminasi variabel dan menyatakan himpunan
penyelesaian sebagai interval. Tanda "#" berarti "tidak sama dengan".
Fungsi ini tidak mendukung fungsi trigonometri (lihat baris cos(x) \textless{} 1/2).
Jika hasil tampak aneh pada pertidaksamaan pecahan atau nilai mutlak,
periksa dengan mensubstitusi beberapa titik uji.

(Bersambung ke Bagian 3: Bahasa Matriks, Submatriks, Pengurutan, Aljabar
Linear, Matriks Simbolik, dan Contoh Bunga/Cicilan.)
* EMT untuk Perhitungan Aljabar (Bagian 3)

Bagian ini mencakup: Bahasa Matriks, Fungsi Pembangun Matriks, Vektorisasi,
Submatriks dan Elemen Matriks, serta Pengurutan dan Pengacakan.

Catatan: keluaran yang berasal dari bilangan acak (random, normal, shuffle)
tidak dicantumkan karena nilainya berbeda setiap dijalankan. Keluaran
perintah latihan baru juga tidak dicantumkan; nilai yang diharapkan ada di
komentar tiap soal.

\begin{eulercomment}
\eulerheading{Bahasa Matriks}
\begin{eulercomment}
Uraian materi: matriks adalah susunan bilangan berbentuk baris dan kolom.
Untuk matriks A (m x n) dan B (n x p), hasil kali AB adalah matriks m x p
dengan unsur

\end{eulercomment}
\begin{eulerformula}
\[
(AB)_{ij}=\sum_{k=1}^{n} a_{ik}b_{kj}
\]
\end{eulerformula}
\begin{eulercomment}
Invers matriks 2x2:

\end{eulercomment}
\begin{eulerformula}
\[
A=\begin{pmatrix}a&b\\c&d\end{pmatrix},\quad A^{-1}=\frac{1}{ad-bc}\begin{pmatrix}d&-b\\-c&a\end{pmatrix}
\]
\end{eulerformula}
\begin{eulercomment}
Dokumentasi inti EMT memuat pembahasan rinci bahasa matriks Euler. Vektor
dan matriks dimasukkan dengan kurung siku, elemen dipisahkan koma, dan baris
dipisahkan titik koma.
\end{eulercomment}
\begin{eulerprompt}
>A=[1,2;3,4]
\end{eulerprompt}
\begin{euleroutput}
              1             2 
              3             4 
\end{euleroutput}
\begin{eulercomment}
Perkalian matriks dinyatakan dengan titik.
\end{eulercomment}
\begin{eulerprompt}
>b=[3;4]
\end{eulerprompt}
\begin{euleroutput}
              3 
              4 
\end{euleroutput}
\begin{eulerprompt}
>b' // transpos b
\end{eulerprompt}
\begin{euleroutput}
  [3,  4]
\end{euleroutput}
\begin{eulerprompt}
>inv(A) //invers A
\end{eulerprompt}
\begin{euleroutput}
             -2             1 
            1.5          -0.5 
\end{euleroutput}
\begin{eulerprompt}
>A.b //perkalian matriks
\end{eulerprompt}
\begin{euleroutput}
             11 
             25 
\end{euleroutput}
\begin{eulerprompt}
>A.inv(A)
\end{eulerprompt}
\begin{euleroutput}
              1             0 
              0             1 
\end{euleroutput}
\begin{eulercomment}
Inti bahasa matriks: semua fungsi dan operator bekerja elemen demi elemen.
\end{eulercomment}
\begin{eulerprompt}
>A.A
\end{eulerprompt}
\begin{euleroutput}
              7            10 
             15            22 
\end{euleroutput}
\begin{eulerprompt}
>A^2 //perpangkatan elemen2 A
\end{eulerprompt}
\begin{euleroutput}
              1             4 
              9            16 
\end{euleroutput}
\begin{eulerprompt}
>A.A.A
\end{eulerprompt}
\begin{euleroutput}
             37            54 
             81           118 
\end{euleroutput}
\begin{eulerprompt}
>power(A,3) //perpangkatan matriks
\end{eulerprompt}
\begin{euleroutput}
             37            54 
             81           118 
\end{euleroutput}
\begin{eulerprompt}
>A/A //pembagian elemen-elemen matriks yang seletak
\end{eulerprompt}
\begin{euleroutput}
              1             1 
              1             1 
\end{euleroutput}
\begin{eulerprompt}
>A/b //pembagian elemen2 A oleh elemen2 b kolom demi kolom (karena b vektor kolom)
\end{eulerprompt}
\begin{euleroutput}
       0.333333      0.666667 
           0.75             1 
\end{euleroutput}
\begin{eulerprompt}
>A\(\backslash\)b // hasilkali invers A dan b, A^(-1)b
\end{eulerprompt}
\begin{euleroutput}
             -2 
            2.5 
\end{euleroutput}
\begin{eulerprompt}
>inv(A).b
\end{eulerprompt}
\begin{euleroutput}
             -2 
            2.5 
\end{euleroutput}
\begin{eulerprompt}
>A\(\backslash\)A   //A^(-1)A
\end{eulerprompt}
\begin{euleroutput}
              1             0 
              0             1 
\end{euleroutput}
\begin{eulerprompt}
>inv(A).A
\end{eulerprompt}
\begin{euleroutput}
              1             0 
              0             1 
\end{euleroutput}
\begin{eulerprompt}
>A*A //perkalian elemen-elemen matriks seletak
\end{eulerprompt}
\begin{euleroutput}
              1             4 
              9            16 
\end{euleroutput}
\begin{eulercomment}
Ini bukan perkalian matriks, melainkan perkalian elemen demi elemen. Hal yang
sama berlaku untuk vektor.
\end{eulercomment}
\begin{eulerprompt}
>b^2 // perpangkatan elemen-elemen matriks/vektor
\end{eulerprompt}
\begin{euleroutput}
              9 
             16 
\end{euleroutput}
\begin{eulercomment}
Jika salah satu operan berupa vektor atau skalar, operan itu diperluas
secara wajar.
\end{eulercomment}
\begin{eulerprompt}
>2*A
\end{eulerprompt}
\begin{euleroutput}
              2             4 
              6             8 
\end{euleroutput}
\begin{eulercomment}
Misalnya jika operan berupa vektor kolom, elemennya diterapkan ke semua baris A.
\end{eulercomment}
\begin{eulerprompt}
>[1,2]*A
\end{eulerprompt}
\begin{euleroutput}
              1             4 
              3             8 
\end{euleroutput}
\begin{eulercomment}
Jika berupa vektor baris, diterapkan ke semua kolom A.
\end{eulercomment}
\begin{eulerprompt}
>A*[2,3]
\end{eulerprompt}
\begin{euleroutput}
              2             6 
              6            12 
\end{euleroutput}
\begin{eulercomment}
Perkalian ini dapat dibayangkan seolah-olah vektor baris v digandakan
hingga berukuran sama dengan A.
\end{eulercomment}
\begin{eulerprompt}
>dup([1,2],2) // dup: menduplikasi/menggandakan vektor [1,2] sebanyak 2 kali (baris)
\end{eulerprompt}
\begin{euleroutput}
              1             2 
              1             2 
\end{euleroutput}
\begin{eulerprompt}
>A*dup([1,2],2)
\end{eulerprompt}
\begin{euleroutput}
              1             4 
              3             8 
\end{euleroutput}
\begin{eulercomment}
Hal ini juga berlaku untuk dua vektor, satu vektor baris dan satu vektor
kolom. Hitung i*j untuk i,j dari 1 sampai 5. Triknya: kalikan 1:5 dengan
transposnya. Bahasa matriks Euler otomatis membuat tabel nilai.
\end{eulercomment}
\begin{eulerprompt}
>(1:5)*(1:5)' // hasilkali elemen-elemen vektor baris dan vektor kolom
\end{eulerprompt}
\begin{euleroutput}
              1             2             3             4             5 
              2             4             6             8            10 
              3             6             9            12            15 
              4             8            12            16            20 
              5            10            15            20            25 
\end{euleroutput}
\begin{eulercomment}
Ingat lagi, ini bukan perkalian matriks!
\end{eulercomment}
\begin{eulerprompt}
>(1:5).(1:5)' // hasilkali vektor baris dan vektor kolom
\end{eulerprompt}
\begin{euleroutput}
  55
\end{euleroutput}
\begin{eulerprompt}
>sum((1:5)*(1:5)) // sama hasilnya
\end{eulerprompt}
\begin{euleroutput}
  55
\end{euleroutput}
\begin{eulercomment}
Operator seperti \textless{} atau == juga bekerja dengan cara yang sama.
\end{eulercomment}
\begin{eulerprompt}
>(1:10)<6 // menguji elemen-elemen yang kurang dari 6
\end{eulerprompt}
\begin{euleroutput}
  [1,  1,  1,  1,  1,  0,  0,  0,  0,  0]
\end{euleroutput}
\begin{eulercomment}
Misalnya, kita dapat menghitung banyak elemen yang memenuhi syarat tertentu
dengan sum().
\end{eulercomment}
\begin{eulerprompt}
>sum((1:10)<6) // banyak elemen yang kurang dari 6
\end{eulerprompt}
\begin{euleroutput}
  5
\end{euleroutput}
\begin{eulercomment}
Euler punya operator pembanding seperti "==" yang memeriksa kesamaan.
Hasilnya vektor 0 dan 1, dengan 1 berarti benar.
\end{eulercomment}
\begin{eulerprompt}
>t=(1:10)^2; t==25 //menguji elemen2 t yang sama dengan 25 (hanya ada 1)
\end{eulerprompt}
\begin{euleroutput}
  [0,  0,  0,  0,  1,  0,  0,  0,  0,  0]
\end{euleroutput}
\begin{eulercomment}
Dari vektor seperti itu, "nonzeros" memilih elemen tak nol. Di sini kita
mendapat indeks semua elemen yang lebih besar dari 50.
\end{eulercomment}
\begin{eulerprompt}
>nonzeros(t>50) //indeks elemen2 t yang lebih besar daripada 50
\end{eulerprompt}
\begin{euleroutput}
  [8,  9,  10]
\end{euleroutput}
\begin{eulercomment}
Vektor indeks ini dapat dipakai untuk mengambil nilai yang bersesuaian di t.
\end{eulercomment}
\begin{eulerprompt}
>t[nonzeros(t>50)] //elemen2 t yang lebih besar daripada 50
\end{eulerprompt}
\begin{euleroutput}
  [64,  81,  100]
\end{euleroutput}
\begin{eulercomment}
Sebagai contoh, cari semua kuadrat dari bilangan 1 sampai 1000 yang bersisa
5 modulo 11 dan bersisa 3 modulo 13.
\end{eulercomment}
\begin{eulerprompt}
>t=1:1000; nonzeros(mod(t^2,11)==5 && mod(t^2,13)==3)
\end{eulerprompt}
\begin{euleroutput}
  [4,  48,  95,  139,  147,  191,  238,  282,  290,  334,  381,  425,
  433,  477,  524,  568,  576,  620,  667,  711,  719,  763,  810,  854,
  862,  906,  953,  997]
\end{euleroutput}
\begin{eulercomment}
EMT kurang efektif untuk perhitungan bilangan bulat karena memakai floating
point presisi ganda secara internal, tetapi sering sangat berguna.

Kita dapat memeriksa kebilangan primaan. Cari berapa banyak kuadrat ditambah 1
yang prima.
\end{eulercomment}
\begin{eulerprompt}
>t=1:1000; length(nonzeros(isprime(t^2+1)))
\end{eulerprompt}
\begin{euleroutput}
  112
\end{euleroutput}
\begin{eulercomment}
Fungsi nonzeros() hanya bekerja untuk vektor. Untuk matriks ada mnonzeros().
\end{eulercomment}
\begin{eulerprompt}
>seed(2); A=random(3,4)
\end{eulerprompt}
\begin{euleroutput}
       0.765761      0.401188      0.406347      0.267829 
        0.13673      0.390567      0.495975      0.952814 
       0.548138      0.006085      0.444255      0.539246 
\end{euleroutput}
\begin{eulercomment}
Fungsi ini mengembalikan indeks elemen yang bukan nol.
\end{eulercomment}
\begin{eulerprompt}
>k=mnonzeros(A<0.4) //indeks elemen2 A yang kurang dari 0,4
\end{eulerprompt}
\begin{euleroutput}
              1             4 
              2             1 
              2             2 
              3             2 
\end{euleroutput}
\begin{eulercomment}
Indeks ini dapat dipakai untuk mengubah elemen menjadi nilai tertentu.
\end{eulercomment}
\begin{eulerprompt}
>mset(A,k,0) //mengganti elemen2 suatu matriks pada indeks tertentu
\end{eulerprompt}
\begin{euleroutput}
       0.765761      0.401188      0.406347             0 
              0             0      0.495975      0.952814 
       0.548138             0      0.444255      0.539246 
\end{euleroutput}
\begin{eulercomment}
mset() juga dapat mengisi elemen pada indeks tersebut dengan elemen matriks
lain.
\end{eulercomment}
\begin{eulerprompt}
>mset(A,k,-random(size(A)))
\end{eulerprompt}
\begin{euleroutput}
       0.765761      0.401188      0.406347     -0.126917 
      -0.122404     -0.691673      0.495975      0.952814 
       0.548138     -0.483902      0.444255      0.539246 
\end{euleroutput}
\begin{eulercomment}
Elemen-elemen juga dapat diambil dalam bentuk vektor.
\end{eulercomment}
\begin{eulerprompt}
>mget(A,k)
\end{eulerprompt}
\begin{euleroutput}
  [0.267829,  0.13673,  0.390567,  0.006085]
\end{euleroutput}
\begin{eulercomment}
Fungsi lain yang berguna adalah extrema, yang mengembalikan nilai minimum
dan maksimum tiap baris beserta posisinya.
\end{eulercomment}
\begin{eulerprompt}
>ex=extrema(A)
\end{eulerprompt}
\begin{euleroutput}
       0.267829             4      0.765761             1 
        0.13673             1      0.952814             4 
       0.006085             2      0.548138             1 
\end{euleroutput}
\begin{eulercomment}
Kita dapat mengambil nilai maksimum tiap baris.
\end{eulercomment}
\begin{eulerprompt}
>ex[,3]'
\end{eulerprompt}
\begin{euleroutput}
  [0.765761,  0.952814,  0.548138]
\end{euleroutput}
\begin{eulercomment}
Hasilnya sama dengan fungsi max().
\end{eulercomment}
\begin{eulerprompt}
>max(A)'
\end{eulerprompt}
\begin{euleroutput}
  [0.765761,  0.952814,  0.548138]
\end{euleroutput}
\begin{eulercomment}
Dengan mget(), kita dapat mengambil indeks dan memakainya untuk mengambil
elemen pada posisi yang sama dari matriks lain.
\end{eulercomment}
\begin{eulerprompt}
>j=(1:rows(A))'|ex[,4], mget(-A,j)
\end{eulerprompt}
\begin{euleroutput}
              1             1 
              2             4 
              3             1 
  [-0.765761,  -0.952814,  -0.548138]
\end{euleroutput}
\begin{eulercomment}
*** Latihan 13.1
Untuk A=[2,1;1,3], hitung det(A), inv(A), dan A.A.
Hasil yang diharapkan: det=5; inv=[0.6,-0.2;-0.2,0.4]; A.A=[5,5;5,10].
\end{eulercomment}
\begin{eulerprompt}
>A=[2,1;1,3]; det(A), inv(A), A.A
\end{eulerprompt}
\begin{euleroutput}
  5
            0.6          -0.2 
           -0.2           0.4 
              5             5 
              5            10 
\end{euleroutput}
\begin{eulercomment}
*** Latihan 13.2
Selesaikan sistem 2x+y=3 dan x+3y=5 dengan operator \textbackslash{} .
Hasil yang diharapkan: x=0.8, y=1.4.
\end{eulercomment}
\begin{eulerprompt}
>A=[2,1;1,3]; A\(\backslash\)[3;5]
\end{eulerprompt}
\begin{euleroutput}
            0.8 
            1.4 
\end{euleroutput}
\begin{eulercomment}
*** Latihan 13.3
Bandingkan A*A (elemen demi elemen) dengan A.A (perkalian matriks).
Hasil yang diharapkan: A*A=[4,1;1,9], sedangkan A.A=[5,5;5,10].
\end{eulercomment}
\begin{eulerprompt}
>A*A, A.A
\end{eulerprompt}
\begin{euleroutput}
              4             1 
              1             9 
              5             5 
              5            10 
\end{euleroutput}
\begin{eulercomment}
*** Latihan 13.4
Buat tabel perkalian 1 sampai 4 dan hitung jumlah diagonal utamanya.
Hasil yang diharapkan: 1+4+9+16 = 30.
\end{eulercomment}
\begin{eulerprompt}
>T=(1:4)*(1:4)'; T, sum(getdiag(T,0))
\end{eulerprompt}
\begin{euleroutput}
              1             2             3             4 
              2             4             6             8 
              3             6             9            12 
              4             8            12            16 
  30
\end{euleroutput}
\begin{eulercomment}
*** Latihan 13.5
Berapa banyak t dari 1 sampai 100 sehingga t\textasciicircum{}2 bersisa 2 modulo 7?
Hasil yang diharapkan: 28 (yaitu t bersisa 3 atau 4 modulo 7).
\end{eulercomment}
\begin{eulerprompt}
>t=1:100; sum(mod(t^2,7)==2)
\end{eulerprompt}
\begin{euleroutput}
  28
\end{euleroutput}
\begin{eulercomment}
Penjelasan: operator "." adalah perkalian matriks, sedangkan "*" dan "\textasciicircum{}"
bekerja elemen demi elemen. Operator "\textbackslash{}" menyelesaikan sistem linear,
kira-kira sama dengan inv(A).b tetapi lebih stabil secara numerik.

\begin{eulercomment}
\eulerheading{Fungsi Pembangun Matriks}
\begin{eulercomment}
Uraian materi: matriks khusus yang sering muncul adalah matriks identitas I
(dengan AI=IA=A), matriks nol, matriks diagonal, dan matriks segitiga.

Untuk membangun matriks, kita dapat menumpuk satu matriks di atas matriks
lain. Jika banyak kolomnya berbeda, yang lebih pendek diisi 0.
\end{eulercomment}
\begin{eulerprompt}
>v=1:3; v_v
\end{eulerprompt}
\begin{euleroutput}
              1             2             3 
              1             2             3 
\end{euleroutput}
\begin{eulercomment}
Demikian pula kita dapat menempelkan matriks di sisi lain, jika banyak
barisnya sama.
\end{eulercomment}
\begin{eulerprompt}
>A=random(3,4); A|v'
\end{eulerprompt}
\begin{euleroutput}
       0.032444     0.0534171      0.595713      0.564454             1 
        0.83916      0.175552      0.396988       0.83514             2 
      0.0257573      0.658585      0.629832      0.770895             3 
\end{euleroutput}
\begin{eulercomment}
Jika banyak barisnya berbeda, matriks yang lebih pendek diisi 0. Ada
pengecualian: bilangan real yang ditempelkan pada matriks dipakai sebagai
kolom yang seluruhnya berisi bilangan itu.
\end{eulercomment}
\begin{eulerprompt}
>A|1
\end{eulerprompt}
\begin{euleroutput}
       0.032444     0.0534171      0.595713      0.564454             1 
        0.83916      0.175552      0.396988       0.83514             1 
      0.0257573      0.658585      0.629832      0.770895             1 
\end{euleroutput}
\begin{eulercomment}
Matriks dapat dibuat dari vektor baris dan kolom.
\end{eulercomment}
\begin{eulerprompt}
>[v;v]
\end{eulerprompt}
\begin{euleroutput}
              1             2             3 
              1             2             3 
\end{euleroutput}
\begin{eulerprompt}
>[v',v']
\end{eulerprompt}
\begin{euleroutput}
              1             1 
              2             2 
              3             3 
\end{euleroutput}
\begin{eulercomment}
Kegunaan utamanya: menafsirkan vektor ekspresi sebagai vektor kolom.
\end{eulercomment}
\begin{eulerprompt}
>"[x,x^2]"(v')
\end{eulerprompt}
\begin{euleroutput}
              1             1 
              2             4 
              3             9 
\end{euleroutput}
\begin{eulercomment}
Untuk mengetahui ukuran A, gunakan fungsi berikut.
\end{eulercomment}
\begin{eulerprompt}
>C=zeros(2,4); rows(C), cols(C), size(C), length(C)
\end{eulerprompt}
\begin{euleroutput}
  2
  4
  [2,  4]
  4
\end{euleroutput}
\begin{eulercomment}
Untuk vektor, ada length().
\end{eulercomment}
\begin{eulerprompt}
>length(2:10)
\end{eulerprompt}
\begin{euleroutput}
  9
\end{euleroutput}
\begin{eulercomment}
Ada banyak fungsi lain yang membangkitkan matriks.
\end{eulercomment}
\begin{eulerprompt}
>ones(2,2)
\end{eulerprompt}
\begin{euleroutput}
              1             1 
              1             1 
\end{euleroutput}
\begin{eulercomment}
Fungsi ini juga dapat dipakai dengan satu parameter. Untuk vektor berisi
bilangan selain 1, gunakan perintah berikut.
\end{eulercomment}
\begin{eulerprompt}
>ones(5)*6
\end{eulerprompt}
\begin{euleroutput}
  [6,  6,  6,  6,  6]
\end{euleroutput}
\begin{eulercomment}
Matriks bilangan acak dapat dibuat dengan random (distribusi seragam) atau
normal (distribusi Gauss).

Fungsi berikut menyusun ulang elemen matriks menjadi matriks lain.
\end{eulercomment}
\begin{eulerprompt}
>redim(1:9,3,3) // menyusun elemen2 1, 2, 3, ..., 9 ke bentuk matriks 3x3
\end{eulerprompt}
\begin{euleroutput}
              1             2             3 
              4             5             6 
              7             8             9 
\end{euleroutput}
\begin{eulercomment}
Dengan fungsi ini dan dup, kita dapat menulis fungsi rep() yang mengulang
vektor sebanyak n kali.
\end{eulercomment}
\begin{eulerprompt}
>function rep(v,n) := redim(dup(v,n),1,n*cols(v))
>rep(1:3,5)
\end{eulerprompt}
\begin{euleroutput}
  [1,  2,  3,  1,  2,  3,  1,  2,  3,  1,  2,  3,  1,  2,  3]
\end{euleroutput}
\begin{eulercomment}
Fungsi multdup() menggandakan elemen vektor.
\end{eulercomment}
\begin{eulerprompt}
>multdup(1:3,5), multdup(1:3,[2,3,2])
\end{eulerprompt}
\begin{euleroutput}
  [1,  1,  1,  1,  1,  2,  2,  2,  2,  2,  3,  3,  3,  3,  3]
  [1,  1,  2,  2,  2,  3,  3]
\end{euleroutput}
\begin{eulercomment}
Fungsi flipx() dan flipy() membalik urutan baris atau kolom matriks.
flipx() membalik secara horizontal.
\end{eulercomment}
\begin{eulerprompt}
>flipx(1:5) //membalik elemen2 vektor baris
\end{eulerprompt}
\begin{euleroutput}
  [5,  4,  3,  2,  1]
\end{euleroutput}
\begin{eulercomment}
Untuk rotasi, Euler punya rotleft() dan rotright().
\end{eulercomment}
\begin{eulerprompt}
>rotleft(1:5) // memutar elemen2 vektor baris
\end{eulerprompt}
\begin{euleroutput}
  [2,  3,  4,  5,  1]
\end{euleroutput}
\begin{eulercomment}
Fungsi khusus drop(v,i) menghapus elemen berindeks i dari vektor v.
\end{eulercomment}
\begin{eulerprompt}
>drop(10:20,3)
\end{eulerprompt}
\begin{euleroutput}
  [10,  11,  13,  14,  15,  16,  17,  18,  19,  20]
\end{euleroutput}
\begin{eulercomment}
Perhatikan bahwa vektor i pada drop(v,i) merujuk ke indeks elemen v, bukan
nilainya. Untuk menghapus berdasarkan nilai, cari dulu indeksnya. Fungsi
indexof(v,x) dapat dipakai untuk mencari elemen x pada vektor terurut v.
\end{eulercomment}
\begin{eulerprompt}
>v=primes(50), i=indexof(v,10:20), drop(v,i)
\end{eulerprompt}
\begin{euleroutput}
  [2,  3,  5,  7,  11,  13,  17,  19,  23,  29,  31,  37,  41,  43,  47]
  [0,  5,  0,  6,  0,  0,  0,  7,  0,  8,  0]
  [2,  3,  5,  7,  23,  29,  31,  37,  41,  43,  47]
\end{euleroutput}
\begin{eulercomment}
Tidak masalah menyertakan indeks di luar jangkauan (seperti 0), indeks
ganda, atau indeks tak terurut.
\end{eulercomment}
\begin{eulerprompt}
>drop(1:10,shuffle([0,0,5,5,7,12,12]))
\end{eulerprompt}
\begin{euleroutput}
  [1,  2,  3,  4,  6,  8,  9,  10]
\end{euleroutput}
\begin{eulercomment}
Ada fungsi khusus untuk mengatur diagonal atau membuat matriks diagonal.
Mulai dari matriks identitas.
\end{eulercomment}
\begin{eulerprompt}
>A=id(5) // matriks identitas 5x5
\end{eulerprompt}
\begin{euleroutput}
              1             0             0             0             0 
              0             1             0             0             0 
              0             0             1             0             0 
              0             0             0             1             0 
              0             0             0             0             1 
\end{euleroutput}
\begin{eulercomment}
Lalu isi diagonal bawah (-1) dengan 1:4.
\end{eulercomment}
\begin{eulerprompt}
>setdiag(A,-1,1:4) //mengganti diagonal di bawah diagonal utama
\end{eulerprompt}
\begin{euleroutput}
              1             0             0             0             0 
              1             1             0             0             0 
              0             2             1             0             0 
              0             0             3             1             0 
              0             0             0             4             1 
\end{euleroutput}
\begin{eulercomment}
Perhatikan bahwa A tidak berubah; kita mendapat matriks baru sebagai hasil
setdiag().

Berikut fungsi yang menghasilkan matriks tridiagonal.
\end{eulercomment}
\begin{eulerprompt}
>function tridiag (n,a,b,c) := setdiag(setdiag(b*id(n),1,c),-1,a); ...
>tridiag(5,1,2,3)
\end{eulerprompt}
\begin{euleroutput}
              2             3             0             0             0 
              1             2             3             0             0 
              0             1             2             3             0 
              0             0             1             2             3 
              0             0             0             1             2 
\end{euleroutput}
\begin{eulercomment}
Diagonal matriks juga dapat diambil. Susun ulang vektor 1:9 menjadi
matriks 3x3.
\end{eulercomment}
\begin{eulerprompt}
>A=redim(1:9,3,3)
\end{eulerprompt}
\begin{euleroutput}
              1             2             3 
              4             5             6 
              7             8             9 
\end{euleroutput}
\begin{eulerprompt}
>d=getdiag(A,0)
\end{eulerprompt}
\begin{euleroutput}
  [1,  5,  9]
\end{euleroutput}
\begin{eulercomment}
Misalnya membagi matriks dengan diagonalnya. Bahasa matriks menerapkan
vektor kolom d pada matriks baris demi baris.
\end{eulercomment}
\begin{eulerprompt}
>fraction A/d'
\end{eulerprompt}
\begin{euleroutput}
          1         2         3 
        4/5         1       6/5 
        7/9       8/9         1 
\end{euleroutput}
\begin{eulercomment}
*** Latihan 14.1
Buat matriks 3x3 berisi 1..9 lalu hitung jejaknya (jumlah diagonal utama).
Hasil yang diharapkan: 1+5+9 = 15.
\end{eulercomment}
\begin{eulerprompt}
>M=redim(1:9,3,3); sum(getdiag(M,0))
\end{eulerprompt}
\begin{euleroutput}
  15
\end{euleroutput}
\begin{eulercomment}
*** Latihan 14.2
Hapus elemen berindeks 2 dan 4 dari vektor 1:10.
Hasil yang diharapkan: [1,3,5,6,7,8,9,10].
\end{eulercomment}
\begin{eulerprompt}
>drop(1:10,[2,4])
\end{eulerprompt}
\begin{euleroutput}
  [1,  3,  5,  6,  7,  8,  9,  10]
\end{euleroutput}
\begin{eulercomment}
*** Latihan 14.3
Buat matriks tridiagonal 4x4 dengan diagonal utama 2 dan diagonal samping 1.
\end{eulercomment}
\begin{eulerprompt}
>tridiag(4,1,2,1)
\end{eulerprompt}
\begin{euleroutput}
              2             1             0             0 
              1             2             1             0 
              0             1             2             1 
              0             0             1             2 
\end{euleroutput}
\begin{eulercomment}
*** Latihan 14.4
Tumpuk vektor 1:3 di atas vektor 4:6, lalu tempelkan kolom bilangan 0.
\end{eulercomment}
\begin{eulerprompt}
>(1:3)_(4:6)|0
\end{eulerprompt}
\begin{euleroutput}
              1             2             3             0 
              4             5             6             0 
\end{euleroutput}
\begin{eulercomment}
*** Latihan 14.5
Balik urutan vektor 1:6, lalu putar ke kiri satu langkah.
Hasil yang diharapkan: [6,5,4,3,2,1] lalu [5,4,3,2,1,6].
\end{eulercomment}
\begin{eulerprompt}
>flipx(1:6), rotleft(flipx(1:6))
\end{eulerprompt}
\begin{euleroutput}
  [6,  5,  4,  3,  2,  1]
  [5,  4,  3,  2,  1,  6]
\end{euleroutput}
\begin{eulercomment}
Penjelasan: setdiag() dan drop() tidak mengubah argumen aslinya, melainkan
mengembalikan matriks/vektor baru. Operator "\textbar{}" menempel ke samping dan "\_"
menumpuk ke bawah.

\begin{eulercomment}
\eulerheading{Vektorisasi}
\begin{eulercomment}
Uraian materi: hampir semua fungsi di Euler bekerja pada vektor dan matriks
jika masuk akal. Misalnya fungsi sqrt() menghitung akar kuadrat semua
elemen. Dengan demikian kita mudah membuat tabel nilai fungsi.
\end{eulercomment}
\begin{eulerprompt}
>sqrt(1:3)
\end{eulerprompt}
\begin{euleroutput}
  [1,  1.41421,  1.73205]
\end{euleroutput}
\begin{eulerprompt}
>x=1:0.01:5; y=log(x)/x^2; // terlalu panjang untuk ditampilkan
\end{eulerprompt}
\begin{eulercomment}
Dengan operator titik dua a:delta:b kita mudah membuat vektor nilai fungsi.
Pada contoh berikut dibuat vektor nilai t[i] dengan jarak 0,1 dari -1 sampai
1, lalu vektor nilai fungsi

\end{eulercomment}
\begin{eulerformula}
\[
s = t^3-t
\]
\end{eulerformula}
\begin{eulerprompt}
>t=-1:0.1:1; s=t^3-t
\end{eulerprompt}
\begin{euleroutput}
  [0,  0.171,  0.288,  0.357,  0.384,  0.375,  0.336,  0.273,  0.192,
  0.099,  0,  -0.099,  -0.192,  -0.273,  -0.336,  -0.375,  -0.384,
  -0.357,  -0.288,  -0.171,  0]
\end{euleroutput}
\begin{eulercomment}
EMT memperluas operator untuk skalar, vektor, dan matriks secara wajar.
Misalnya vektor kolom dikali vektor baris diperluas menjadi matriks. Pada
perintah berikut v' adalah vektor transpos (vektor kolom).
\end{eulercomment}
\begin{eulerprompt}
>shortest (1:5)*(1:5)'
\end{eulerprompt}
\begin{euleroutput}
       1      2      3      4      5 
       2      4      6      8     10 
       3      6      9     12     15 
       4      8     12     16     20 
       5     10     15     20     25 
\end{euleroutput}
\begin{eulercomment}
Perhatikan, ini berbeda dengan perkalian matriks yang dinyatakan dengan
titik "." di EMT.
\end{eulercomment}
\begin{eulerprompt}
>(1:5).(1:5)'
\end{eulerprompt}
\begin{euleroutput}
  55
\end{euleroutput}
\begin{eulercomment}
Secara default, vektor baris dicetak dalam format ringkas.
\end{eulercomment}
\begin{eulerprompt}
>[1,2,3,4]
\end{eulerprompt}
\begin{euleroutput}
  [1,  2,  3,  4]
\end{euleroutput}
\begin{eulercomment}
Untuk matriks, operator . menyatakan perkalian matriks dan A' menyatakan
transpos. Matriks 1x1 dapat dipakai seperti bilangan real.
\end{eulercomment}
\begin{eulerprompt}
>v:=[1,2]; v.v', %^2
\end{eulerprompt}
\begin{euleroutput}
  5
  25
\end{euleroutput}
\begin{eulerprompt}
>v=1:4; v'
\end{eulerprompt}
\begin{euleroutput}
              1 
              2 
              3 
              4 
\end{euleroutput}
\begin{eulercomment}
Kita dapat menghitung matriks A kali vektor b.
\end{eulercomment}
\begin{eulerprompt}
>A=[1,2,3,4;5,6,7,8]; A.v'
\end{eulerprompt}
\begin{euleroutput}
             30 
             70 
\end{euleroutput}
\begin{eulercomment}
Perhatikan v tetap vektor baris. Jadi v'.v berbeda dengan v.v'.
\end{eulercomment}
\begin{eulerprompt}
>v'.v
\end{eulerprompt}
\begin{euleroutput}
              1             2             3             4 
              2             4             6             8 
              3             6             9            12 
              4             8            12            16 
\end{euleroutput}
\begin{eulercomment}
v.v' menghitung kuadrat norma v untuk vektor baris v. Hasilnya vektor 1x1
yang bekerja seperti bilangan real.
\end{eulercomment}
\begin{eulerprompt}
>v.v'
\end{eulerprompt}
\begin{euleroutput}
  30
\end{euleroutput}
\begin{eulerprompt}
>norm(v)^2
\end{eulerprompt}
\begin{euleroutput}
  30
\end{euleroutput}
\begin{eulercomment}
Ringkasan aturannya:

- Fungsi yang diterapkan pada vektor atau matriks diterapkan ke tiap elemen.
- Operator pada dua matriks berukuran sama diterapkan elemen demi elemen.
- Jika dimensinya berbeda, keduanya diperluas secara wajar sampai berukuran sama.

Misalnya nilai skalar dikali vektor mengalikan nilai itu ke tiap elemen.
\end{eulercomment}
\begin{eulerprompt}
>[1,2,3]^2
\end{eulerprompt}
\begin{euleroutput}
  [1,  4,  9]
\end{euleroutput}
\begin{eulerprompt}
>v:=[1,2,3]; v*v'
\end{eulerprompt}
\begin{euleroutput}
              1             2             3 
              2             4             6 
              3             6             9 
\end{euleroutput}
\begin{eulercomment}
Hasil kali skalar memakai perkalian matriks, bukan *!
\end{eulercomment}
\begin{eulerprompt}
>v.v'
\end{eulerprompt}
\begin{euleroutput}
  14
\end{euleroutput}
\begin{eulercomment}
Ada banyak fungsi untuk matriks. Berikut daftar singkatnya; baca dokumentasi
untuk informasi lebih lanjut.

\end{eulercomment}
\begin{eulerttcomment}
  sum,prod menghitung jumlah dan hasil kali tiap baris
  cumsum,cumprod melakukan hal yang sama secara kumulatif
  max,min menghitung nilai ekstrem tiap baris
  extrema mengembalikan vektor informasi ekstrem
  diag(A,i) mengembalikan diagonal ke-i
  setdiag(A,i,v) mengatur diagonal ke-i
  id(n) matriks identitas
  det(A) determinan
  charpoly(A) polinomial karakteristik
  eigenvalues(A) nilai eigen
\end{eulerttcomment}
\begin{eulerprompt}
>v*v, sum(v*v), cumsum(v*v)
\end{eulerprompt}
\begin{euleroutput}
  [1,  4,  9]
  14
  [1,  5,  14]
\end{euleroutput}
\begin{eulercomment}
Operator : membangkitkan vektor baris berjarak sama, opsional dengan langkah.
\end{eulercomment}
\begin{eulerprompt}
>1:4, 1:2:10
\end{eulerprompt}
\begin{euleroutput}
  [1,  2,  3,  4]
  [1,  3,  5,  7,  9]
\end{euleroutput}
\begin{eulercomment}
Untuk menggabung matriks dan vektor ada operator "\textbar{}" dan "\_".
\end{eulercomment}
\begin{eulerprompt}
>[1,2,3]|[4,5], [1,2,3]_1
\end{eulerprompt}
\begin{euleroutput}
  [1,  2,  3,  4,  5]
              1             2             3 
              1             1             1 
\end{euleroutput}
\begin{eulercomment}
Elemen matriks dirujuk dengan "A[i,j]".
\end{eulercomment}
\begin{eulerprompt}
>A:=[1,2,3;4,5,6;7,8,9]; A[2,3]
\end{eulerprompt}
\begin{euleroutput}
  6
\end{euleroutput}
\begin{eulercomment}
Untuk vektor baris atau kolom, v[i] adalah elemen ke-i. Untuk matriks,
A[i] mengembalikan seluruh baris ke-i.
\end{eulercomment}
\begin{eulerprompt}
>v:=[2,4,6,8]; v[3], A[3]
\end{eulerprompt}
\begin{euleroutput}
  6
  [7,  8,  9]
\end{euleroutput}
\begin{eulercomment}
Indeks juga dapat berupa vektor indeks. ":" berarti semua indeks.
\end{eulercomment}
\begin{eulerprompt}
>v[1:2], A[:,2]
\end{eulerprompt}
\begin{euleroutput}
  [2,  4]
              2 
              5 
              8 
\end{euleroutput}
\begin{eulercomment}
Bentuk singkat ":" adalah menghilangkan indeks sama sekali.
\end{eulercomment}
\begin{eulerprompt}
>A[,2:3]
\end{eulerprompt}
\begin{euleroutput}
              2             3 
              5             6 
              8             9 
\end{euleroutput}
\begin{eulercomment}
Untuk keperluan vektorisasi, elemen matriks dapat diakses seperti vektor.
\end{eulercomment}
\begin{eulerprompt}
>A\{4\}
\end{eulerprompt}
\begin{euleroutput}
  4
\end{euleroutput}
\begin{eulercomment}
Matriks juga dapat diratakan dengan redim(), yang diimplementasikan dalam
fungsi flatten().
\end{eulercomment}
\begin{eulerprompt}
>redim(A,1,prod(size(A))), flatten(A)
\end{eulerprompt}
\begin{euleroutput}
  [1,  2,  3,  4,  5,  6,  7,  8,  9]
  [1,  2,  3,  4,  5,  6,  7,  8,  9]
\end{euleroutput}
\begin{eulercomment}
Untuk memakai matriks sebagai tabel, reset ke format default lalu hitung
tabel nilai sinus dan cosinus. Sudut default dalam radian.
\end{eulercomment}
\begin{eulerprompt}
>defformat; w=0°:45°:360°; w=w'; deg(w)
\end{eulerprompt}
\begin{euleroutput}
  Variable  not found!
  Error in degree
  Error in:
  defformat; w=0°:45°:360°; w=w'; deg(w) ...
                     ^
\end{euleroutput}
\begin{eulercomment}
Sekarang tambahkan kolom ke matriks.
\end{eulercomment}
\begin{eulerprompt}
>M = deg(w)|w|cos(w)|sin(w)
\end{eulerprompt}
\begin{euleroutput}
  Wrong argument.
  Cannot use a string matrix here.
  
  Try "trace errors" to inspect local variables after errors.
  deg:
      return x/pi*180;
  Error in:
  M = deg(w)|w|cos(w)|sin(w) ...
            ^
\end{euleroutput}
\begin{eulercomment}
Dengan bahasa matriks kita dapat membuat beberapa tabel beberapa fungsi
sekaligus. Misalnya hitung t[j]\textasciicircum{}i untuk i dari 1 sampai n. Kita mendapat
matriks dengan setiap barisnya adalah tabel t\textasciicircum{}i untuk satu i, yaitu unsur

\end{eulercomment}
\begin{eulerformula}
\[
a_{i,j} = t_j^i, \quad 1 \le j \le 101, \quad 1 \le i \le n
\]
\end{eulerformula}
\begin{eulercomment}
Fungsi yang tidak bekerja untuk masukan vektor harus "divektorkan" dengan
kata kunci "map" pada definisi fungsi. Fungsi akan dievaluasi untuk setiap
elemen parameter vektor.

Integrasi numerik integrate() hanya bekerja untuk batas interval skalar,
sehingga perlu divektorkan.
\end{eulercomment}
\begin{eulerprompt}
>function map f(x) := integrate("x^x",1,x)
>f([1:5])
\end{eulerprompt}
\begin{euleroutput}
  [0,  2.05045,  13.7251,  113.336,  1241.03]
\end{euleroutput}
\begin{eulercomment}
*** Latihan 15.1
Hitung t\textasciicircum{}2 untuk t=0:0.5:2. Hasil yang diharapkan: [0,0.25,1,2.25,4].
\end{eulercomment}
\begin{eulerprompt}
>t=0:0.5:2; t^2
\end{eulerprompt}
\begin{euleroutput}
  [0,  0.25,  1,  2.25,  4]
\end{euleroutput}
\begin{eulercomment}
*** Latihan 15.2
Hitung 1+2+...+100 dan 5!. Hasil yang diharapkan: 5050 dan 120.
\end{eulercomment}
\begin{eulerprompt}
>sum(1:100), prod(1:5)
\end{eulerprompt}
\begin{euleroutput}
  5050
  120
\end{euleroutput}
\begin{eulercomment}
*** Latihan 15.3
Hitung jumlah kumulatif 1:5. Hasil yang diharapkan: [1,3,6,10,15].
\end{eulercomment}
\begin{eulerprompt}
>cumsum(1:5)
\end{eulerprompt}
\begin{euleroutput}
  [1,  3,  6,  10,  15]
\end{euleroutput}
\begin{eulercomment}
*** Latihan 15.4
Hitung deret 1 + 1/2 + 1/4 + ... + 1/2\textasciicircum{}10 (jumlah 11 suku).
Hasil yang diharapkan: 2047/1024 (mendekati 2).
\end{eulercomment}
\begin{eulerprompt}
>fraction sum(1/2^(0:10))
\end{eulerprompt}
\begin{euleroutput}
  2047/1024
\end{euleroutput}
\begin{eulercomment}
*** Latihan 15.5
Buat tabel n, n\textasciicircum{}2, n\textasciicircum{}3 untuk n=1:5.
\end{eulercomment}
\begin{eulerprompt}
>n=(1:5)'; n|n^2|n^3
\end{eulerprompt}
\begin{euleroutput}
              1             1             1 
              2             4             8 
              3             9            27 
              4            16            64 
              5            25           125 
\end{euleroutput}
\begin{eulercomment}
Penjelasan: dengan notasi a:b, kita mengganti perulangan dengan satu
ekspresi vektor. Menulis sum(1/2\textasciicircum{}(0:10)) setara dengan loop yang
menjumlahkan 11 suku.

\begin{eulercomment}
\eulerheading{Submatriks dan Elemen Matriks}
\begin{eulercomment}
Uraian materi: submatriks adalah matriks yang diambil dari sebagian baris
dan kolom matriks lain.

Untuk mengakses elemen matriks, gunakan notasi kurung siku.
\end{eulercomment}
\begin{eulerprompt}
>A=[1,2,3;4,5,6;7,8,9], A[2,2]
\end{eulerprompt}
\begin{euleroutput}
              1             2             3 
              4             5             6 
              7             8             9 
  5
\end{euleroutput}
\begin{eulercomment}
Kita dapat mengakses satu baris penuh.
\end{eulercomment}
\begin{eulerprompt}
>A[2]
\end{eulerprompt}
\begin{euleroutput}
  [4,  5,  6]
\end{euleroutput}
\begin{eulercomment}
Pada vektor baris atau kolom, ini mengembalikan satu elemen.
\end{eulercomment}
\begin{eulerprompt}
>v=1:3; v[2]
\end{eulerprompt}
\begin{euleroutput}
  2
\end{euleroutput}
\begin{eulercomment}
Agar mendapat baris pertama pada matriks 1xn maupun mxn, tentukan semua
kolom dengan indeks kedua kosong.
\end{eulercomment}
\begin{eulerprompt}
>A[2,]
\end{eulerprompt}
\begin{euleroutput}
  [4,  5,  6]
\end{euleroutput}
\begin{eulercomment}
Jika indeks berupa vektor indeks, Euler mengembalikan baris-baris yang
bersesuaian. Di sini kita ingin baris pertama dan kedua dari A.
\end{eulercomment}
\begin{eulerprompt}
>A[[1,2]]
\end{eulerprompt}
\begin{euleroutput}
              1             2             3 
              4             5             6 
\end{euleroutput}
\begin{eulercomment}
Kita bahkan dapat menyusun ulang A dengan vektor indeks. (A sendiri tidak
berubah; yang dihitung adalah versi yang disusun ulang.)
\end{eulercomment}
\begin{eulerprompt}
>A[[3,2,1]]
\end{eulerprompt}
\begin{euleroutput}
              7             8             9 
              4             5             6 
              1             2             3 
\end{euleroutput}
\begin{eulercomment}
Trik indeks juga bekerja untuk kolom. Contoh ini memilih semua baris A dan
kolom kedua dan ketiga.
\end{eulercomment}
\begin{eulerprompt}
>A[1:3,2:3]
\end{eulerprompt}
\begin{euleroutput}
              2             3 
              5             6 
              8             9 
\end{euleroutput}
\begin{eulercomment}
Untuk singkatnya, ":" berarti semua indeks baris atau kolom.
\end{eulercomment}
\begin{eulerprompt}
>A[:,3]
\end{eulerprompt}
\begin{euleroutput}
              3 
              6 
              9 
\end{euleroutput}
\begin{eulercomment}
Atau kosongkan indeks pertama.
\end{eulercomment}
\begin{eulerprompt}
>A[,2:3]
\end{eulerprompt}
\begin{euleroutput}
              2             3 
              5             6 
              8             9 
\end{euleroutput}
\begin{eulercomment}
Kita juga dapat mengambil baris terakhir A.
\end{eulercomment}
\begin{eulerprompt}
>A[-1]
\end{eulerprompt}
\begin{euleroutput}
  [7,  8,  9]
\end{euleroutput}
\begin{eulercomment}
Sekarang ubah elemen A dengan memberi nilai pada submatriks. Ini benar-benar
mengubah matriks A yang tersimpan.
\end{eulercomment}
\begin{eulerprompt}
>A[1,1]=4
\end{eulerprompt}
\begin{euleroutput}
              4             2             3 
              4             5             6 
              7             8             9 
\end{euleroutput}
\begin{eulercomment}
Kita juga dapat memberi nilai pada satu baris A.
\end{eulercomment}
\begin{eulerprompt}
>A[1]=[-1,-1,-1]
\end{eulerprompt}
\begin{euleroutput}
             -1            -1            -1 
              4             5             6 
              7             8             9 
\end{euleroutput}
\begin{eulercomment}
Bahkan pada submatriks, asalkan ukurannya sesuai.
\end{eulercomment}
\begin{eulerprompt}
>A[1:2,1:2]=[5,6;7,8]
\end{eulerprompt}
\begin{euleroutput}
              5             6            -1 
              7             8             6 
              7             8             9 
\end{euleroutput}
\begin{eulercomment}
Beberapa pintasan juga diperbolehkan.
\end{eulercomment}
\begin{eulerprompt}
>A[1:2,1:2]=0
\end{eulerprompt}
\begin{euleroutput}
              0             0            -1 
              0             0             6 
              7             8             9 
\end{euleroutput}
\begin{eulercomment}
Peringatan: indeks di luar batas mengembalikan matriks kosong atau pesan
kesalahan, bergantung pada pengaturan sistem. Defaultnya pesan kesalahan.
Ingat juga bahwa indeks negatif dapat dipakai untuk mengakses elemen dari
belakang.
\end{eulercomment}
\begin{eulerprompt}
>A[4]
\end{eulerprompt}
\begin{euleroutput}
  Row index 4 out of bounds!
  Error in:
  A[4] ...
      ^
\end{euleroutput}
\begin{eulercomment}
*** Latihan 16.1
Untuk A=[1,2,3;4,5,6;7,8,9], ambil elemen baris 2 kolom 3, kolom kedua,
dan baris ketiga.
Hasil yang diharapkan: 6; [2;5;8]; [7,8,9].
\end{eulercomment}
\begin{eulerprompt}
>A=[1,2,3;4,5,6;7,8,9]; A[2,3], A[:,2], A[3]
\end{eulerprompt}
\begin{euleroutput}
  6
              2 
              5 
              8 
  [7,  8,  9]
\end{euleroutput}
\begin{eulercomment}
*** Latihan 16.2
Ambil submatriks sudut: baris dan kolom 1 dan 3.
Hasil yang diharapkan: [1,3;7,9].
\end{eulercomment}
\begin{eulerprompt}
>A[[1,3],[1,3]]
\end{eulerprompt}
\begin{euleroutput}
              1             3 
              7             9 
\end{euleroutput}
\begin{eulercomment}
*** Latihan 16.3
Ubah elemen tengah A[2,2] menjadi 0, lalu hitung determinan A. (Karena A
diubah, jalankan ulang definisi A sebelum latihan berikutnya.)
\end{eulercomment}
\begin{eulerprompt}
>A[2,2]=0; det(A)
\end{eulerprompt}
\begin{euleroutput}
  60
\end{euleroutput}
\begin{eulercomment}
*** Latihan 16.4
Tukar baris pertama dan ketiga dari matriks 3x3 berisi 1..9.
Hasil yang diharapkan: baris [7,8,9], [4,5,6], [1,2,3].
\end{eulercomment}
\begin{eulerprompt}
>M=redim(1:9,3,3); M[[3,2,1]]
\end{eulerprompt}
\begin{euleroutput}
              7             8             9 
              4             5             6 
              1             2             3 
\end{euleroutput}
\begin{eulercomment}
*** Latihan 16.5
Jumlahkan elemen-elemen submatriks baris 1:2 dan kolom 2:3 dari M.
Hasil yang diharapkan: 2+3+5+6 = 16.
\end{eulercomment}
\begin{eulerprompt}
>sum(sum(M[1:2,2:3]))
\end{eulerprompt}
\begin{euleroutput}
              5 
             11 
\end{euleroutput}
\begin{eulercomment}
Penjelasan: indeks vektor memilih beberapa baris/kolom sekaligus. Memberi
nilai pada A[i,j] mengubah matriks tersimpan, sedangkan A[[3,2,1]] hanya
menghitung salinan yang disusun ulang.

\begin{eulercomment}
\eulerheading{Pengurutan dan Pengacakan}
\begin{eulercomment}
Uraian materi: pengurutan menyusun elemen secara menaik. Indeks hasil
pengurutan berguna untuk menyusun ulang vektor lain dengan urutan yang sama.

Fungsi sort() mengurutkan vektor baris.
\end{eulercomment}
\begin{eulerprompt}
>sort([5,6,4,8,1,9])
\end{eulerprompt}
\begin{euleroutput}
  [1,  4,  5,  6,  8,  9]
\end{euleroutput}
\begin{eulercomment}
Sering kita perlu mengetahui indeks vektor terurut pada vektor asal. Ini
dapat dipakai untuk menyusun ulang vektor lain dengan urutan yang sama.

Mari kita acak sebuah vektor (hasil acak berbeda tiap dijalankan).
\end{eulercomment}
\begin{eulerprompt}
>v=shuffle(1:10)
\end{eulerprompt}
\begin{euleroutput}
  [5,  3,  7,  10,  6,  2,  1,  4,  9,  8]
\end{euleroutput}
\begin{eulercomment}
Indeks memuat urutan yang benar dari v.
\end{eulercomment}
\begin{eulerprompt}
>\{vs,ind\}=sort(v); v[ind]
\end{eulerprompt}
\begin{euleroutput}
  [1,  2,  3,  4,  5,  6,  7,  8,  9,  10]
\end{euleroutput}
\begin{eulercomment}
Ini juga bekerja untuk vektor string.
\end{eulercomment}
\begin{eulerprompt}
>s=["a","d","e","a","aa","e"]
\end{eulerprompt}
\begin{euleroutput}
  a
  d
  e
  a
  aa
  e
\end{euleroutput}
\begin{eulerprompt}
>\{ss,ind\}=sort(s); ss
\end{eulerprompt}
\begin{euleroutput}
  a
  a
  aa
  d
  e
  e
\end{euleroutput}
\begin{eulercomment}
Seperti terlihat, posisi entri ganda agak acak.
\end{eulercomment}
\begin{eulerprompt}
>ind
\end{eulerprompt}
\begin{euleroutput}
  [4,  1,  5,  2,  6,  3]
\end{euleroutput}
\begin{eulercomment}
Fungsi unique mengembalikan daftar terurut dari elemen-elemen unik vektor.
\end{eulercomment}
\begin{eulerprompt}
>intrandom(1,10,10), unique(%)
\end{eulerprompt}
\begin{euleroutput}
  [2,  3,  8,  1,  4,  4,  9,  2,  6,  5]
  [1,  2,  3,  4,  5,  6,  8,  9]
\end{euleroutput}
\begin{eulercomment}
Ini juga bekerja untuk vektor string.
\end{eulercomment}
\begin{eulerprompt}
>unique(s)
\end{eulerprompt}
\begin{euleroutput}
  a
  aa
  d
  e
\end{euleroutput}
\begin{eulercomment}
*** Latihan 17.1
Urutkan [3,1,2,5,4] secara menaik lalu menurun.
Hasil yang diharapkan: [1,2,3,4,5] dan [5,4,3,2,1].
\end{eulercomment}
\begin{eulerprompt}
>sort([3,1,2,5,4]), flipx(sort([3,1,2,5,4]))
\end{eulerprompt}
\begin{euleroutput}
  [1,  2,  3,  4,  5]
  [5,  4,  3,  2,  1]
\end{euleroutput}
\begin{eulercomment}
*** Latihan 17.2
Dapatkan unsur unik dari [3,1,3,2,1]. Hasil yang diharapkan: [1,2,3].
\end{eulercomment}
\begin{eulerprompt}
>unique([3,1,3,2,1])
\end{eulerprompt}
\begin{euleroutput}
  [1,  2,  3]
\end{euleroutput}
\begin{eulercomment}
*** Latihan 17.3
Urutkan vektor nama berdasarkan nilai mereka.
Nilai [70,90,80] untuk nama ["Ani","Budi","Citra"] menghasilkan urutan
Ani, Citra, Budi.
\end{eulercomment}
\begin{eulerprompt}
>nilai=[70,90,80]; nama=["Ani","Budi","Citra"]; \{ns,idx\}=sort(nilai); nama[idx]
\end{eulerprompt}
\begin{euleroutput}
  Ani
  Citra
  Budi
\end{euleroutput}
\begin{eulercomment}
*** Latihan 17.4
Tentukan median dari [7,3,9,1,5]. Hasil yang diharapkan: 5 (elemen tengah
setelah diurutkan: 1,3,5,7,9).
\end{eulercomment}
\begin{eulerprompt}
>sort([7,3,9,1,5])[3]
\end{eulerprompt}
\begin{euleroutput}
  [7,  3,  9,  1,  5]
\end{euleroutput}
\begin{eulercomment}
Penjelasan: sort() dengan dua keluaran \{nilai\_terurut, indeks\} memungkinkan
mengurutkan vektor lain dengan urutan yang sama. unique() sekaligus
menghapus duplikat dan mengurutkan.

(Bersambung ke Bagian 4: Aljabar Linear, Matriks Simbolik, Nilai Numerik
dalam Ekspresi Simbolik, dan Contoh Bunga/Cicilan.)
\end{eulercomment}
\end{eulernotebook}
\end{document}
